\documentclass{article}
\usepackage[preprint]{neurips_2026}

\usepackage[utf8]{inputenc}
\usepackage[T1]{fontenc}
\usepackage{hyperref}
\usepackage{url}
\usepackage{booktabs}
\usepackage{amsfonts}
\usepackage{amsmath}
\usepackage{nicefrac}
\usepackage{microtype}
\usepackage{graphicx}
\usepackage{subcaption}
\usepackage{xcolor}

\graphicspath{{figs/}}

% A last-resort loosening pass, so that a paragraph carrying a long unbreakable
% token (\texttt{} operator names, "leave-one-operator-out", inline intervals)
% ends up with a slightly loose line instead of one that runs into the margin.
\setlength{\emergencystretch}{3em}

\definecolor{cexact}{HTML}{0072B2}
\definecolor{cblock}{HTML}{D55E00}
\definecolor{crouter}{HTML}{009E73}

\title{How the Dirty Set Moves: A Benchmark for\\Delta-Input Sparse Matrix
Multiplication}

\author{%
  Anonymous Authors\\
  \texttt{}\\
  % Submission version is anonymized. The camera-ready author block is held
  % outside this artifact on purpose -- it ships publicly, so naming the authors
  % here would de-anonymize the submission. Restore it at camera-ready from
  % docs/PAPER_AUTHORS.md in the umbrella repo.
}

% Float placement. The stock limits leave at most one float per page and a
% 20% text minimum, which in a paper with this float density pushed three
% tables past their first reference and onto the last two pages.
\renewcommand{\topfraction}{0.9}
\renewcommand{\bottomfraction}{0.7}
\renewcommand{\textfraction}{0.08}
\renewcommand{\floatpagefraction}{0.8}
\setcounter{topnumber}{3}
\setcounter{bottomnumber}{2}
\setcounter{totalnumber}{4}

% Placeholder for a number whose measurement is still running. Renders in bold
% brackets so it cannot survive a proof-read, and `make check` greps for it.
\newcommand{\PENDING}[1]{\textbf{[PENDING: #1]}}


\begin{document}
\maketitle

\begin{abstract}
A large class of numerical workloads applies a \emph{fixed} sparse operator $A$
to a state that changes by a sparse increment each step, so the product can be
maintained rather than recomputed: $A(x{+}dx) = Ax + A\,dx$. Power-flow solvers,
circuit simulators, and implicit FEM codes all have this shape. Established
sparse suites time a single product against a static operand and therefore
cannot see it. We define delta-input SpMM as a benchmarkable task and supply a
benchmark: 21 SuiteSparse operators, a dirty column set of fixed size that moves
between steps, a churn rate $\rho$, five seeds, and two \emph{motion models}
differing only in where the replacement columns come from --- structural
neighbours of the survivors under \emph{local drift}, a fresh uniform draw under
\emph{uniform jump} --- with a mixing parameter making them endpoints of one axis.

\textbf{The motion model, not the operator's structure or the churn rate alone,
decides what the benchmark measures.} Under local drift a block-scheduled delta
is the faster arm in 418 of 418 churning cells; under uniform jump, at the same
operators, the same $|D|$, the same $\rho$ and the same seeds, it is faster in
313 of 419 and arm selection becomes a real problem --- though a second,
independently written block-scheduled implementation wins 395 of 419 there, so
part of that selection problem belongs to one implementation rather than to the
family. The same switch changes what
a global row merge is worth ($2.87\times$ against $1.74\times$ at $\rho \geq
0.01$) and the size of the gap between the two implementation families. A
control separates the two localities confounded in that contrast: under a random
symmetric relabelling, which preserves the graph and the structural trajectory
and destroys index locality, the block-scheduled arm wins 306--386 of 418 at
$1.97$--$2.37\times$ rather than $5.62\times$ over four permutations. Index
locality is doing real work, a median $2.39\times$ $[1.93, 2.93]$ excluding $1$
in all four; whether motion locality contributes beyond it is not resolved at 21
operators.

We report four results that a benchmark of this kind needs and that this one
initially lacked. \textbf{A family claim needs more than one member per family},
so we write and release an open block-scheduled reference arm against our own
externally-visible description of the strategy: it reproduces the dichotomy
independently --- losing on a frozen set and winning 418 of 418 churning cells
under local drift, the same count the proprietary arm reaches --- and makes the
merge result --- previously
derived from schedules only a closed implementation builds --- an integer
identity a reader can recompute. \textbf{A motion model needs evidence that
anything moves that way}, so we instrument two solvers and read the dirty set off
solver state rather than generating it: transient propagation lands on the
local-drift endpoint exactly ($\alpha_{\text{emp}} = 0.000$ on all nine
operators, at $\rho_{\text{emp}} \approx 0.008$), while Newton convergence sits
near the middle of the axis at ten times the churn ($\alpha_{\text{emp}} \approx
0.44$, $\rho_{\text{emp}} \approx 0.27$) --- supporting the regime claim for
time-stepping workloads and withdrawing it for iterative solves.
And \textbf{a delta benchmark will silently charge one family for work the other
never does}: we found two such taxes pointing in opposite directions, one an
$M{\times}B$ zero-fill inside session construction that changed the sign of an
aggregate comparison when removed, the other a change-detection cost $29\times$
more expensive on one side than the other for the same question.
Finally, \textbf{two conditions that differ in the intended variable may also
differ in how much work there is}: uniform jump destroys locality and separately
shrinks the dirty slice to $0.415\times$ its starting nonzeros, so a third motion
model that holds the work while destroying locality raises the uniform-jump
margin by $1.58\times$ and makes every drift-versus-jump contrast here an upper
bound.

Where selection does exist, the win belongs to a model-\emph{selection}
procedure rather than to any single model named in advance: fitted offline it
beats an online policy that probes both arms on geomean regret by $0.044$--$0.078$
in all three scopes carrying a live decision, while a pre-registered logistic
regression loses to the policy on two of them. A short cost model explains why
the one structural predictor tried had to fail --- it prices the apply and says
nothing about re-preparation, which is the only half the motion model touches.
\end{abstract}

\section{Introduction}
\label{sec:intro}

Sparse matrix--vector and matrix--matrix products are among the most heavily
benchmarked kernels in scientific computing. The standard measurement is
one-shot: load an operator, apply it once to a dense operand, report bandwidth
or FLOP rate. That protocol assumes the operand is unrelated to the previous
one.

A great many real solvers violate that assumption in the same way. The operator
is a discretization or a network admittance matrix and does not change; the
state does, and it changes locally. A contingency analysis perturbs a handful of
buses. A circuit simulator's Newton step updates the nodes that have not yet
converged. An implicit FEM solve reloads a region under a moving load. In each
case $x_{t+1} = x_t + dx_t$ with $dx_t$ supported on a small set $D_t$ of
columns, and the product can be maintained rather than recomputed
(Figure~\ref{fig:task}).

\begin{figure}[tbp]
  \centering
  \includegraphics[width=\linewidth]{fig_task}
  \caption{The delta-input SpMM task. $A$ is fixed across steps; the state
  changes by a sparse increment supported on a dirty column set $D$. The
  incremental product costs $\mathrm{nnz}(A[:,D])$ rather than
  $\mathrm{nnz}(A)$, but only if the implementation can restrict itself to $D$
  without re-preparing itself whenever $D$ changes.}
  \label{fig:task}
\end{figure}

The arithmetic saving is not what makes this hard. A delta implementation must
\emph{prepare} itself for a particular dirty set, and preparation is expensive
relative to a step. When $D$ is stable, preparation amortizes. When $D$ moves,
preparation may have to be redone, and the implementation that was fastest on a
frozen set can become the slowest. No sparse benchmark we are aware of varies
the dirty set, so none of them can express this trade-off.

Varying it exposes a second variable that is easy to miss. Two workloads can
agree on the operator, on $|D|$, on the churn rate $\rho$, and on the random
seed, and still disagree about which implementation is faster, because they
disagree about \emph{where the replacement columns come from}. We call that
choice the motion model. It decides the size of the opportunity, whether an
implementation-selection problem exists at all, and what a particular schedule
optimization is worth.

\paragraph{Contributions.}
\begin{enumerate}
\itemsep2pt
\item \textbf{A task definition and benchmark} (\S\ref{sec:bench}). Delta-input
SpMM with a moving dirty set: 21 SuiteSparse operators from power flow, circuit
simulation and structural FEM; two motion models and a mixing parameter that
interpolates between them; five churn rates including $\rho{=}0.002$; five
dirty-set seeds; an open and deliberately strong baseline; and a five-part
metric protocol (\S\ref{sec:metrics}) that makes amortization assumptions
impossible to leave implicit.
\item \textbf{The motion model separates the space} (\S\ref{sec:families}).
Under local drift the block-scheduled arm is faster in 418 of 418 churning
cells and the alternative class is empty. Under uniform jump it is faster in 313
of 419 and both classes are populated. The operator set, $|D|$, $\rho$ and the
seeds are identical; only the motion differs.
\item \textbf{Motion locality and index locality, separated}
(\S\ref{sec:control}). A control that destroys index locality while holding the
graph and the structural trajectory fixed, replicated over four independent
permutations, costs the margin a median $2.39\times$ $[1.93, 2.93]$ and 32--45
of 418 cells, and $1.97$--$2.05\times$ at a bound that
over-corrects for the baseline's own relabelling penalty; destroying motion
locality costs $2.92\times$ $[2.43, 3.49]$ and 105 cells. Ranking the two
against each other is a $1.02$--$1.23\times$ contrast whose interval covers $1$
in every permutation, and we report it as unresolved rather than as a result. The mixing parameter
turns the dichotomy into a curve, and most of the fall has happened by
$\alpha{=}0.25$.
\item \textbf{What a merged schedule is worth is a property of the motion}
(\S\ref{sec:merge}), measured on 21 operators and five seeds, with two
counter-examples reported.
\item \textbf{Selection, and what it needs to observe} (\S\ref{sec:select},
Appendix~\ref{sec:amortize}). Under local drift there is nothing to select. Under
uniform jump there is, and an offline model-selection procedure beats an online
policy that probes both arms; no single pre-registered model does.
\item \textbf{A negative result with a named mechanism} (Appendix~\ref{sec:patch}), and
\textbf{a measurement-methodology result} (\S\ref{sec:defect}).
\end{enumerate}

\section{Background and related work}
\label{sec:related}

\paragraph{Sparse kernel benchmarks.}
The SuiteSparse Matrix Collection~\citep{davis2011} is the standard operator
source, and the surrounding benchmark literature --- SpMV performance
studies~\citep{williams2007, vuduc2003}, format autotuning~\citep{im2004},
graph-kernel suites~\citep{beamer2015}, and vendor suites --- measures the
one-shot product. The operand is arbitrary and its relationship to the previous
operand is not part of the workload. These suites cannot distinguish an
implementation that is fast per step from one that is fast per step \emph{and}
cheap to re-prepare, which is the distinction that decides the problem studied
here.

\paragraph{Sparse compilers.}
Sparse tensor compilers generate a kernel from an index expression and a chosen
format~\citep{kjolstad2017, ye2023}, and a scheduling language decides the loop
order and blocking. The compiled artifact of Section~\ref{sec:families} is a
degenerate instance: one compiled segment per column-tile is a scheduling
decision that these systems could express. What they do not model is the cost of
\emph{re}-compiling when the sparsity of the operand changes, because their
compilation is indexed by the operator's format and not by which columns are
live this step. That is the axis this benchmark varies, and it is orthogonal to
the code-generation quality those systems optimize --- a point worth stating
plainly, because the two are easy to conflate: a better generated kernel helps
both families here equally.

\paragraph{Incremental and streaming computation.}
Incremental view maintenance~\citep{gupta1995} and its higher-order delta
treatment in DBToaster~\citep{koch2014}, self-adjusting
computation~\citep{acar2002, acar2005}, differential
dataflow~\citep{mcsherry2013}, and dynamic graph
algorithms~\citep{demetrescu2004} all treat the same underlying idea ---
recompute only what changed. Self-adjusting computation and differential
dataflow are the closest in spirit, since both maintain a dependence structure
and repair it when inputs move; the difference is that they repair a general
dataflow and pay a per-change bookkeeping cost, where a delta schedule is a flat
array over one fixed operator and the whole question is what re-preparing that
array costs. Streaming graph containers --- STINGER~\citep{ediger2012},
Aspen~\citep{dhulipala2019}, GraphOne~\citep{kumar2020},
Terrace~\citep{pandey2021} --- do vary an update stream, but the thing that
changes there is the \emph{operator}, and the batch is scored on ingestion
throughput. Here the operator is fixed and the operand moves, which is the
opposite decomposition. GraphBLAS~\citep{davis2019, buluc2017, graphblas2023}
exposes masked operations that can express a restriction to $D$. We measured
it rather than asserting anything about it, and the result is in
\S\ref{sec:baselines}: the delta form is not the object of study there, and it
shows. We are
not aware of a benchmark in any of these communities that varies the stability
of the changed set, let alone the manner of its movement.

\paragraph{Updating a factorization instead of a product.}
The classical answer to ``the operand changed a little'' is not to redo the
product at all but to update a factorization or an inverse by a low-rank
correction~\citep{sherman1950, hager1989, golub2013}. That is a different task
and it is worth saying why it is not a baseline here. Sherman--Morrison and its
descendants update $A^{-1}$ or a factorization of $A$ when \emph{$A$} changes by
a low-rank term; the workload studied here holds $A$ fixed and changes $x$. The
product $A\,dx$ with $dx$ sparse is already the cheapest correct form, so there
is no factorization to amortize --- the whole cost is deciding which entries of
$A$ to touch and how to reach them.

\paragraph{Kernel selection and autotuning.}
Selecting a sparse format or kernel from matrix features is a well-developed
line of work, usually as offline supervised learning~\citep{li2013,
sedaghati2015, zhao2018}, and usually evaluated by classification accuracy with
random train/test splits. Section~\ref{sec:select} argues that both of those
choices are wrong for this problem: random splits leak operator identity, and
accuracy is the wrong loss when the penalty for a mistake spans an order of
magnitude and a large share of the labels sits inside the measurement floor.
Online alternatives --- bandit-style selection and profile-guided dispatch ---
are the natural comparison, and the policy we evaluate is a simple instance of
that family.

\section{The benchmark}
\label{sec:bench}

\subsection{Task}

Fix a sparse $A \in \mathbb{R}^{N\times M}$. At each step $t$ the caller
supplies a dirty column set $D_t \subseteq \{1,\dots,M\}$ and a dense block
$dx_t \in \mathbb{R}^{|D_t| \times B}$, and requires
\[
  dy_t \;=\; A[:,D_t]\,dx_t \;\in\; \mathbb{R}^{N\times B},
\]
exactly, in floating point, to the accuracy a direct evaluation would give. $B$
is the batch width. An implementation may perform arbitrary preparation, may
cache it across steps, and is charged for all of it. Sequences are long: the
regime of interest is thousands of steps against one operator, which is what
makes preparation worth doing and also what makes stale preparation worth
avoiding.

\subsection{Motion models}
\label{sec:motion}

The dirty set is the workload. We generate $D_0$ by breadth-first search from a
random seed column over the operator's own adjacency, giving a connected
neighbourhood of a requested size --- the structural counterpart of a local
disturbance. At each subsequent step a fraction $\rho$ of $D_t$ is retired.
$|D|$ is held constant, so the same number of columns is added back, and the
motion model is the rule that chooses them (Figure~\ref{fig:workload}):

\begin{itemize}
\itemsep1pt
\item \textbf{local drift} refills from the structural neighbourhood of the
survivors, so the set migrates through the topology;
\item \textbf{uniform jump} redraws uniformly from every live column, so the set
scatters;
\item \textbf{mix($\alpha$)} draws a fraction $\alpha$ of each step's
replacements uniformly and the rest from the neighbourhood, so local drift and
uniform jump are $\alpha{=}0$ and $\alpha{=}1$ of one axis. Both endpoints
delegate to the same functions with the same arguments and therefore the same
draws, so a mix run at $\alpha \in \{0,1\}$ reproduces the corresponding drift
or jump run trajectory for trajectory rather than in distribution.
\end{itemize}

We sweep $\rho \in \{0, 0.002, 0.01, 0.05, 0.25\}$ over five dirty-set seeds.
$\rho{=}0$ is the frozen set that informal measurements of this idea implicitly
assume; it is a special case rather than the default, and it is the identity
for both motion models, which makes the two $\rho{=}0$ blocks the same workload
measured twice and therefore a usable noise floor (\S\ref{sec:limits}).

\begin{figure}[tbp]
  \centering
  \includegraphics[width=\linewidth]{fig_workload}
  \caption{The two motion models on one fixed topology, at each churn rate.
  Both retire the same fraction of the dirty set per step and hold $|D|$
  constant; they differ only in where the replacements come from. Fixed
  topology is what makes the operator reusable, and the disturbed region is
  what a delta implementation must track. They are not the same thing, and
  \S\ref{sec:families}--\S\ref{sec:select} are organized around the gap.}
  \label{fig:workload}
\end{figure}

$\rho$ is a rate and it quantizes. Three of the 1050 measured cells carry a
label $\rho > 0$ but never move their dirty set, because $\lceil \rho|D|
\rceil$ rounds to zero or the retired column is immediately re-added; all three
are one operator at one seed whose BFS returns 138 columns where every other
seed returns 256. They are frozen cells wearing a churn label and are excluded
from the churning population, with the excluded keys reported rather than
absorbed. That leaves 418 churning cells under local drift and 419 under uniform
jump, against 105 frozen cells each.

Two controls bracket the model and are reported alongside it: \texttt{contiguous}
(a best case for any blocked layout) and \texttt{scattered} (uniformly random
columns, a worst case). Neither is realistic; both exist so that a submission
cannot be tuned to the neighbourhood generator without the effect being visible.

\paragraph{A third model, measured and withdrawn.}
We also measured a \texttt{cyclic} model that rotates through a small ring of
recurring dirty sets. It is not reportable and no number from it appears in this
paper. The ring is 4--5 distinct sets wide; one arm caches eight prepared sets
and the other caches one, so every \texttt{cyclic} ratio restates that pair of
integers rather than anything about the workload. It becomes reportable when the
baseline arm gains a multi-entry cache, and not before.

\subsection{Do real solvers actually drift?}
\label{sec:realmotion}

Everything downstream is conditioned on a motion model, and the argument that
local drift is the realistic one is made in \S\ref{sec:intro} by assertion: a
disturbance propagates through a fixed topology along that topology. Both models
in \S\ref{sec:motion} are generators. \texttt{drift} refills from a BFS
neighbourhood because it was told to and \texttt{jump} redraws uniformly because
it was told to; neither is an observation, and the realism claim is doing enough
work here that it should not rest on one.

We therefore produce dirty sets that nobody chose. Two solvers run on the same
SuiteSparse topologies, and at each step the dirty set is read off the solver
state --- the $|D|$ unknowns that moved most --- rather than generated. Holding
$|D|$ fixed at the benchmark's own 256 is deliberate: a threshold instead makes
the answer an artifact of the threshold, and on a diffusion with a held source
the disturbed region grows monotonically, nothing ever retires, and the measured
churn rate is identically zero.

\texttt{transient} steps an implicit linear diffusion on the operator's own
graph after a pulse disturbance at one node --- the shape a linearized swing
equation, an RC ladder and heat conduction all share, and the ``contingency
perturbs a handful of buses'' workload of \S\ref{sec:intro}. Nothing places the
resulting front; where it goes is whatever solving says. \texttt{newton} runs
Newton--Raphson on a nonlinear system with the same sparsity, where the dirty
set is the set of unknowns not yet converged --- the ``circuit simulator updates
the nodes that have not yet converged'' workload, and a structurally different
mechanism.

Each trajectory is then scored on the paper's own axis: $\rho_{\text{emp}}$, the
fraction of the dirty set that turned over, and $\alpha_{\text{emp}}$, the
fraction of \emph{arriving} columns that were not structural neighbours of the
previous set. $\alpha_{\text{emp}}$ is exactly the mixing parameter of
\S\ref{sec:motion}, measured instead of set.

\paragraph{Result.} Over nine operators drawn from the three domains
\S\ref{sec:intro} names, five seeds and 200 steps each:

\begin{itemize}\itemsep2pt
  \item \texttt{transient}: $\alpha_{\text{emp}} = 0.000$ on \emph{all nine}
    operators --- median and range alike --- at $\rho_{\text{emp}}$ median
    $0.0078$ (range $0.0039$--$0.0156$). Every arriving column is a structural
    neighbour of a surviving one, and the set turns over at four to sixteen
    columns in a thousand per step.
  \item \texttt{newton}: $\alpha_{\text{emp}}$ median $0.439$ (range
    $0.067$--$0.787$) at $\rho_{\text{emp}}$ median $0.2695$ (range
    $0.078$--$0.496$). Nearly half the arrivals are unrelated to the survivors,
    and the set turns over an order of magnitude faster.
\end{itemize}

Transient propagation therefore lands exactly on the local-drift endpoint of the
axis, and at the churn rate the headline regime occupies: $\rho \approx 0.008$
sits between the benchmark's $\rho{=}0.002$ and $\rho{=}0.01$ rows, where
\S\ref{sec:families} measures $5.65$ and $5.39\times$ for the open reference arm. That is the first
measurement behind a claim this paper previously made by assertion, and it is a
stronger result than we expected --- not ``mostly local'' but $\alpha = 0$ on every
operator tried.

Newton convergence does not. At $\alpha \approx 0.44$ and $\rho \approx 0.27$ it
sits near the middle of the mixing axis at its highest churn rate, which is
where \S\ref{sec:mixing} shows most of the block-scheduled advantage has already
decayed. The two solvers do not agree, and that is the finding: transient propagation sits at the local-drift end of the axis at low
churn, which is the regime the headline occupies, while Newton convergence does
not --- its dirty set both churns faster and arrives further from the survivors.
Section~\ref{sec:intro} names power flow, circuit simulation and implicit FEM
together as workloads that drift; on this evidence that is right for the
time-stepping ones and wrong for the iterative-solve one, and the paper's regime
claim should be read as being about the former.

\paragraph{What this is and is not.} These are dirty sets produced by a solve on
a real topology with a real sparsity pattern, rather than by a rule. They are
not traces captured from a production energy-management system, a commercial
circuit simulator, or a deployed FEM code: the forcing is synthetic and the
nonlinearity in \texttt{newton} is chosen rather than measured. What they
establish is the weaker but load-bearing claim that \S\ref{sec:intro} previously
asserted with no evidence at all --- that on a fixed topology, a propagating
disturbance produces a dirty set whose arrivals are overwhelmingly structural
neighbours of its survivors. A trace from a deployed solver remains the right
next measurement, and we do not have one.

\subsection{Operators}

Twenty-one SuiteSparse operators, chosen to span the domains where the task
arises rather than to flatter any implementation, and including shapes expected
to be hostile: power flow (\texttt{TSOPF}$\times 3$, \texttt{case9}), power
networks (\texttt{bcspwr09/10}), circuit simulation (\texttt{scircuit},
\texttt{memplus}, \texttt{add32}, \texttt{circuit\_3}, \texttt{rajat03}), device
simulation (\texttt{wang3}), and structural and shell FEM (\texttt{ct20stif},
\texttt{bcsstk18/25}, \texttt{msc10848}, \texttt{nasa2910}, \texttt{s3rmt3m3},
\texttt{olafu}, \texttt{raefsky3}, \texttt{epb2}). Sizes span $N=1{,}723$ to
$170{,}998$ and $\mathrm{nnz}=6{,}511$ to $2{,}600{,}295$; $|D|$ is 256 except
where BFS returns fewer, down to 138. Table~\ref{tab:operators} lists them with
results.

\subsection{Baselines}
\label{sec:baselines}

A delta kernel can be made to look arbitrarily good by comparing it against a
weak delta. Table~\ref{tab:baselines} is the ladder we measured.

\begin{table}[tbp]
  \centering
  \caption{The baseline ladder: geomean ms/step over the ten configurations
  (five operators $\times$ two orderings) for which all rungs were measured,
  $B{=}8$. The benchmark reports against the \emph{bottom} rung. Choosing a rung
  near the top inflates any delta implementation by one to two orders of
  magnitude, and doing so silently is the most common defect we found in our own
  earlier measurements.}
  \label{tab:baselines}
  \small
  \begin{tabular}{lrr}
\toprule
delta implementation & ms/step & $\div$ our baseline \\
\midrule
PyTorch, delta re-sliced each step & 3.331 & 197.9$\times$ \\
PyTorch \texttt{torch.sparse.mm}, full recompute & 1.234 & 73.3$\times$ \\
PyTorch, delta pre-sliced once & 0.383 & 22.7$\times$ \\
CSC delta, hand-written C++ & 0.059 & 3.5$\times$ \\
column-exact CSR delta, 1 thread & 0.041 & 2.5$\times$ \\
\textbf{column-exact CSR delta, best thread count} & 0.017 & 1.0$\times$ \\
\bottomrule
\end{tabular}

\end{table}

The reference baseline is a \textbf{column-exact CSR delta}: for the current $D$
it compiles a compressed row list holding exactly the nonzeros of $A[:,D]$,
templated on batch width and optionally threaded. It is open source and is not
our product; we invested in making it strong because a benchmark whose baseline
is weak measures nothing. A naive PyTorch delta that re-slices $A[:,D]$ each
step is slower than recomputing the whole product ($3.331$ vs $1.234$~ms): the
obvious way to write this inside a framework is a pessimization.

\paragraph{GraphBLAS, measured.}
GraphBLAS~\citep{davis2019, graphblas2023} is the system in the literature
closest to this task, so we put it on the ladder rather than describing it.
SuiteSparse:GraphBLAS through \texttt{python-graphblas} 2025.2.0, one thread,
on the ladder's own ten configurations, geomean ms/step at $B{=}8$:

\begin{itemize}\itemsep2pt
  \item $A \, dx$ with $dx$ a sparse operand populated only on the dirty rows ---
    the formulation that most directly says ``the operand is a delta'' ---
    $8.04$~ms, $473\times$ our baseline.
  \item $A \, dx$ with a dense operand, the formulation \S\ref{sec:related}
    describes: $3.72$~ms, $219\times$.
  \item Extracting $A[:, D]$ each step and multiplying: $2.31$~ms, $136\times$.
    This is the best of the three that are legal for a moving dirty set.
  \item The same extract hoisted out of the loop --- legal only on a
    \emph{frozen} set, so not comparable with any churning cell here, and
    reported to show what GraphBLAS costs when re-preparation is free:
    $0.38$~ms, $22\times$.
\end{itemize}

Two things are worth separating in that. The mask-on-a-full-traversal reading is
supported --- the dense-operand form is $219\times$ --- but it is not the whole
story, because the formulation that \emph{does} restrict to $D$ is still
$136\times$, and hoisting the restriction away entirely still leaves $22\times$.
The gap is not mainly about mask semantics; it is that a general runtime for sparse linear
algebra carries per-call machinery this task cannot amortize.

The honest caveats: this is measured through the Python binding, so some of the
per-step cost is binding overhead a C caller would not pay, and we did not
attempt to tune descriptors. We report it as evidence that the delta form is not
what this ecosystem is built for, not as a tuned comparison, and every
formulation agrees with a direct evaluation to $3.5 \times 10^{-16}$.

\paragraph{End-to-end comparison.}
For a practitioner writing this inside a framework today, adopting the
benchmark's best configuration is worth $8.8\times$/$15.2\times$/$22.9\times$ at
$B{=}1$/$8$/$16$, geomean over the same ten configurations. Most of that
is leaving the framework rather than choosing between delta implementations: the
open baseline alone is $22.7\times$ faster than the best in-framework delta
(Table~\ref{tab:baselines}). Everything after \S\ref{sec:families} concerns the
residual factor.

Two independently written harnesses measured the hand-written CSC delta on the
same operators and configurations; their timings agree to within $0.1$--$12.6\%$
per configuration (e.g.\ \texttt{TSOPF} under RCM: $1.5740$ vs $1.5752$~ms), which
we take as a weak but real check on the timing methodology.

\subsection{Metrics}
\label{sec:metrics}

Every cell reports all five numbers, including cells where the system under test
loses. The protocol exists because each of the first four can be made to look
good alone.

\begin{enumerate}
\itemsep1pt
\item \textbf{Steady-state cost per step} with preparation amortized.
\item \textbf{One-shot cost}: a single application including all preparation.
\item \textbf{Breakeven step count}: how many steps before (2) is repaid.
\item \textbf{Cost under churn} at each $\rho$ and each motion model, with
re-preparation charged where it falls.
\item \textbf{Maximum relative error} against a direct evaluation.
\end{enumerate}

\paragraph{Aggregation, stated once.}
Ratios are reported as the median across the five dirty-set seeds within an
(operator, $\rho$) cell, then the geometric mean over operators. Per-cell win
counts are reported separately and are counts of cells, not of operators. Local
drift is never averaged into uniform jump anywhere in this paper.

Frozen and churning cells are averaged into one another only in a scope that is
named a \emph{workload mixture} at every appearance, and never silently. Two of
the three scopes that carry a live selection decision (\S\ref{sec:contrast}) are
such mixtures, so this is not a corner case: the mixture is where $\rho$ has any
variance at all, and therefore the only scope in which ``is the churn rate the
decisive variable'' is a well-posed question. It is also the scope in which the
decision is dominated by a trivial gate at $\rho{>}0$, which
\S\ref{sec:select} and \S\ref{sec:contrast} report rather than hide. Every number
below names its regime --- frozen, churning, or the mixture --- its motion
model, and its $n$.

\section{Two implementation families}
\label{sec:families}

Delta implementations divide by \emph{what they prepare for}, and that choice
fixes both their per-step cost and when they must re-prepare. A column-exact
implementation prepares for exactly $D$, so it touches no wasted entries but must
re-prepare whenever $D$ changes at all. A block-scheduled implementation prepares
per block-column, so it scans padding --- a geomean $1.65\times$ more entries
than needed on the ten-configuration ladder --- but survives movement that stays
inside the blocks it already covers.

\paragraph{Three systems, so that a family claim has more than one member.}
Two implementations cannot establish a claim about two families: every
difference would be confounded with every other engineering difference between
one open C++ header and one closed runtime. We therefore measure \emph{three}.

The column-exact arm is our open baseline. The block-scheduled arm is a
proprietary system that we characterize only by its measured cost profile:
per-step cost, re-preparation frequency as a function of $\rho$ and of the
motion model, and re-preparation cost. Its internal data structures are not
disclosed and are not required by anything we claim.

The third is an \textbf{open block-scheduled reference arm}, written for this
paper against the externally-visible strategy stated above and nothing else, and
released with the benchmark. It shares no source with the proprietary arm, and
two of its choices differ on purpose so that a result which reproduces on it is a
property of the strategy rather than of one codebase: its dx scratch is
tile-local rather than $M$-wide, so it cannot incur the session-construction tax
of \S\ref{sec:defect} at all, and its compiled segments are position-independent,
so its preparation cache never invalidates --- the strongest form of the family's
``survives movement'' property, which gives the block-scheduled side the benefit of
the doubt. It is not tuned to win. Its purpose is falsification: if the
dichotomy below did not appear on it, \S\ref{sec:families} would be a statement
about one product and this paper would have to say so.

It also carries \S\ref{sec:merge}, which measures what a globally row-merged
schedule is worth. That section's row counts previously came from schedules only
the closed arm builds, and were the one part of this paper not reproducible from
the artifact at all; the reference arm emits both a per-tile and a merged form,
so the quantity is now an integer identity a reader can recompute.

The benchmark of \S\ref{sec:bench} is defined without reference to any arm, and
the selection policy of \S\ref{sec:select} treats all of them as black boxes.
Section~\ref{sec:limits} states what remains closed and
\S\ref{sec:artifact} describes what we release.

\paragraph{Frozen dirty set.}
On a frozen dirty set the column-exact baseline is the better arm. The
block-scheduled arm reaches $0.49$/$0.58$/$0.59$ of it at $B{=}1$/$8$/$16$ over
the ten-configuration ladder (Figure~\ref{fig:frozen}, left). Over the
21-operator sweep at $B{=}8$ it is the better arm in 3 of 105 frozen cells under
local drift and 5 of 105 under uniform jump, and in each case it is at parity
rather than past it.

Padding largely accounts for that margin and does not close it
(Figure~\ref{fig:frozen}, right). Across the 21 operators the correlation
between padding and frozen slowdown is $+0.94$ (drift) and $+0.93$ (jump), but
the per-operator residual --- slowdown divided by padding --- spans
$0.99$--$1.83\times$. \texttt{raefsky3} pads $1.00\times$ and still loses
$1.38\times$; \texttt{add32} pads $1.52\times$ and loses $2.78\times$. At
$B{=}1$ the ladder's gap is $2.04\times$ against $1.65\times$ padding. So the
block-scheduled arm's best possible frozen result is a tie, reached only where
padding vanishes, and the gap it actually pays is wider than padding alone
predicts by up to $1.8\times$. We did not isolate the remainder.

\begin{figure}[tbp]
  \centering
  \includegraphics[width=\linewidth]{fig_frozen}
  \caption{Frozen dirty set. Left: the ratio at three batch widths over the ten
  configurations for which all batch widths were measured; the thick line is the
  geomean. Right: per-operator frozen slowdown against padding on the
  21-operator sweep, one point per operator per motion model, median over five
  seeds. The dashed line is what padding alone would predict. This is the half
  of the picture a one-shot benchmark can see, and on its own it would retire
  block scheduling as an idea.}
  \label{fig:frozen}
\end{figure}

\paragraph{Moving dirty set.}
Once the set moves the ordering reverses, and by how much depends on the motion
model. At $\rho{=}0.002$ --- two columns in a thousand --- the column-exact arm
must recompile, and recompilation costs tens of steps; the block-scheduled arm
re-prepares only for blocks it has not already covered. Under local drift the
geomean of always-block-scheduled over always-column-exact is $6.94\times$ at
$\rho{=}0.002$ and stays above $4.7\times$ through $\rho{=}0.25$; under uniform
jump it opens at $2.95\times$ and decays to $1.35\times$ (21 operators, five
seeds). Figure~\ref{fig:churn} plots the better arm rather than a fixed one, so
its curves read $6.94 \to 4.74$ and $2.97 \to 1.54$. What the decay measures is
block reuse: uniform redraw destroys the reuse the block-scheduled arm depends
on and local drift preserves it.

The uniform-jump half of that comparison is an \emph{upper bound} on what
destroying locality costs, not an estimate of it. Uniform redraw also shrinks the
dirty slice's nonzero count --- to $0.415\times$ its starting value over 400
steps, against $0.818\times$ under drift --- which helps the column-exact arm
more than the block-scheduled one, whose cost has a floor set by the tiles it
must scan. Holding the work constant and destroying only locality raises the
uniform-jump margin by $1.58\times$ (\S\ref{sec:limits}). Every uniform-jump
figure in this paper is the uncorrected one.

\begin{figure}[tbp]
  \centering
  \begin{subfigure}{0.46\linewidth}
    \includegraphics[width=\linewidth]{fig_churn_response}
    \caption{Margin of the better arm over the baseline against $\rho$.}
    \label{fig:churn}
  \end{subfigure}\hfill
  \begin{subfigure}{0.50\linewidth}
    \includegraphics[width=\linewidth]{fig_heatmap}
    \caption{Which arm wins, and by how much, for every operator $\times$ churn
    rate. Blue is column-exact, orange block-scheduled.}
    \label{fig:heatmap}
  \end{subfigure}
  \caption{The same operators, the same $|D|$, the same $\rho$ and the same
  seeds under both motion models. The local-drift panel of
  Figure~\ref{fig:heatmap} has no blue cell at any $\rho>0$; the uniform-jump
  panel keeps a real mixture.}
\end{figure}

\paragraph{The dichotomy reproduces on an independently written arm.}
Everything above is measured on the proprietary arm, so it establishes a
property of one system. We re-measure with the open reference arm added as a
third lane, under the same generator, the same $|D|$, the same seeds and the
same 400 steps.

\textbf{This is a separate sweep and its numbers are not interchangeable with
the ones above.} It runs five lanes rather than four, rotates lane order by step
rather than holding it fixed, and detects a moved tile set with the $O(|D|)$
check of \S\ref{sec:defect} rather than the one the main sweep used. The
proprietary arm therefore reads differently in it, and both arms are quoted from
it below so that the comparison is within one file. Table~\ref{tab:family} gives
both.

\begin{table}[tbp]
  \centering
  \caption{The three-arm sweep: 21 operators, five seeds, 400 steps, $B{=}8$,
  one thread, churn-labelled cells that never moved excluded. Both
  block-scheduled arms are measured in the same run against the same
  column-exact baseline; the main sweep's figures for the proprietary arm
  (\S\ref{sec:families}, above) come from a different harness and are not
  comparable cell by cell.}
  \label{tab:family}
  \small
  \begin{tabular}{llrrrr}
    \toprule
    & & \multicolumn{2}{c}{\emph{local drift}} & \multicolumn{2}{c}{\emph{uniform jump}} \\
    \cmidrule(lr){3-4}\cmidrule(lr){5-6}
    arm & regime & cells won & geomean & cells won & geomean \\
    \midrule
    open reference & frozen    & 5/105   & $0.37\times$ & 10/105  & $0.38\times$ \\
    proprietary    & frozen    & 0/105   & $0.41\times$ & 3/105   & $0.42\times$ \\
    \addlinespace
    open reference & churning  & \textbf{418/418} & $5.20\times$ & \textbf{395/419} & $2.08\times$ \\
    proprietary    & churning  & 399/418 & $4.45\times$ & 237/419 & $1.39\times$ \\
    \bottomrule
  \end{tabular}
\end{table}

\begin{figure}[tbp]
  \centering
  \includegraphics[width=\linewidth]{fig_family}
  \caption{Two implementations of one strategy against the shared column-exact
  baseline, geomean over 21 operators, median across five seeds, churn-labelled
  cells that never moved excluded. The open reference arm was written against
  this paper's prose and shares no code with the proprietary one. What matters
  is that the curves keep the same shape under local drift: the frozen loss, the
  churning win and the direction of the decay are properties of the strategy.
  Where they separate --- uniform jump, at every churn rate --- is where a claim
  is about a product rather than a family, and it is the scope
  \S\ref{sec:select} operates in.}
  \label{fig:family}
\end{figure}

Read down the table, the shape is the same for both: a frozen loss of about
$0.4\times$, a churning win under local drift on essentially every cell, and a
smaller churning win under uniform jump. Read across, the open arm is ahead
everywhere the set moves. The frozen loss, the churning win, the ordering of the
two motion models and the direction of the decay along each are therefore
properties of the strategy rather than of one codebase, which is what the
control was built to establish (Figure~\ref{fig:family}).

Within this sweep the open arm's local-drift curve runs $5.65 / 5.39 / 4.88 /
4.94\times$ across $\rho \in \{0.002, 0.01, 0.05, 0.25\}$ and the proprietary
arm's runs $5.82 / 4.89 / 3.75 / 3.69\times$: they start together and separate as
churn rises. Under uniform jump they are $2.93 / 2.20 / 1.80 / 1.62\times$
against $2.07 / 1.48 / 1.16 / 1.04\times$.

What does \emph{not} transfer between sweeps is the level. The main sweep's
proprietary arm runs $6.94 \to 4.74\times$ under local drift where the open arm
here runs $5.65 \to 4.94\times$ --- the same $418/418$ win count and the same
high-churn endpoint to within $4\%$, but $19\%$ apart at $\rho{=}0.002$. Under
uniform jump the low-churn ends agree to under $1\%$ ($2.93$ against $2.95\times$)
and the high-churn ends do not ($1.62$ against $1.35\times$). We therefore treat
the win counts as the reproducible part and the geomeans as harness-dependent,
which is the honest split: which arm was faster on a cell survives a change of
harness, and a ratio of means does not.

The two churning counts come from applying \S\ref{sec:motion}'s own exclusion
rule, and it is worth saying how they arrived at it. Of the 420 drift cells with
a churn label, the open arm loses exactly two --- and those two are precisely the
degenerate cells \S\ref{sec:motion} identifies independently, one operator at one
seed whose BFS returns 138 columns rather than 256, where $\lceil \rho|D| \rceil$
rounds to zero and the dirty set never actually moves. The arm loses exactly
where nothing moved, which is what a frozen cell wearing a churn label should do.
That the exclusion set was derived from the generator and the loss set from an
independent implementation's timings, and that they coincide exactly, is a check
on both.

\paragraph{But how much of the \emph{margin} is the strategy?}
Comparing the two block-scheduled arms directly on identical cells, the shipped
arm is faster on a frozen set ($1.12\times$) and \emph{slower} once the set moves
--- $0.86\times$ under local drift and $0.66\times$ under uniform jump, geomean
$0.75\times$ over all churning cells (per-cell range $0.53$--$1.33\times$). The open reference arm was not tuned; it
avoids the $M{\times}B$ scratch of \S\ref{sec:defect} by construction and keeps
its compiled segments position-independent, and those two choices are apparently
worth more under churn than whatever the shipped arm does better.

This matters beyond a ranking, because it moves the scope in which
\S\ref{sec:select}'s question is live. Under uniform jump the shipped arm is the
better one in 237 of 419 churning cells --- both classes populated, a real
selection problem --- while the open arm is the better one in 395 of 419, which
is much closer to the one-sided picture local drift gives. The open arm's 24
uniform-jump losses are not spread evenly either: 16 are \texttt{scircuit}, 5
\texttt{circuit\_3} and 3 \texttt{wang3}, the three lowest-density
circuit-simulation operators, so what survives of the selection problem is
concentrated rather than general. \textbf{So a substantial part
of ``uniform jump makes arm selection a real problem'' is a property of one
implementation rather than of the family.} We report \S\ref{sec:select} as
measured, on the arm pair it was measured with, and flag that a better
block-scheduled implementation shrinks the scope in which any selector is worth
having. That flag is not compiler-specific: on the 378 cells built both ways,
the open arm wins 115 of 126 uniform-jump churning cells against the proprietary
arm's 66 under g++, and 97 against 63 under clang (\S\ref{sec:limits}). This is
the clearest case in the paper of a conclusion that two implementations could
not have reached.

\paragraph{Which arm wins, per cell.}
Per cell rather than per operator, the block-scheduled arm is the faster of the
two in \textbf{418 of 418} churning cells under local drift --- every operator,
every seed, every churn rate measured --- and in 313 of 419 under uniform jump.
Under local drift, at $\rho{=}0.25$, it remains the better arm on all 21
operators; under uniform jump the count falls from 20 of 21 at $\rho{=}0.002$ to
12 of 21 at $\rho{=}0.25$.

Those cells are not near-ties. The \emph{smallest} local-drift churning margin
is $1.123\times$, and \textbf{zero} of the 418 sit within $10\%$ of parity
against 53 of 419 under uniform jump. So under local drift there is no
implementation-selection problem in the measured range and any policy that
switches once the set moves is optimal by construction, while under uniform jump
there is one. Section~\ref{sec:control} asks how much of that separation is the
motion model and how much is the index order these operators happen to ship
in.

\section{Two one-sided measurement taxes}
\label{sec:defect}

Neither defect below is a defect in either implementation. Both are defects in
what the harness charges to which of them, both were found only by asking of
every timed operation whether the other family pays it too, and both survived
several rounds of review. We report them together and early because they are the
most transferable thing here: the benchmark's own protocol (\S\ref{sec:metrics})
is what makes them findable, and a delta benchmark without that discipline will
have them without knowing it.

\paragraph{Session construction, charged to one family.}
Each arm's re-preparation was timed around the call that produces a prepared
session, which is what a cache miss costs the caller. One family of
implementations allocated and zero-filled an $M \times B$ scratch block inside
that call and outside its own build timer; the other compiles a $|D| \times B$
scratch about three orders of magnitude smaller. The zero-fill was therefore
charged to every cell of one family and to no cell of the other. Measured
directly it was $0.0145$--$0.4159$~ms per session on a $0.32$--$10.44$~MB block,
rising with $N$.

Pooling the scratch across sessions removes it. The residual is
$0.0083$--$0.0086$~ms across all five operators, a $1.04\times$ spread over a
$33\times$ range of scratch sizes, which is the test that the fill is gone rather
than relocated. The fix was not to move the timer inside the build, which would
have made the measurement self-consistent while still charging one family for
session construction.

The effect is not confined to noise. On 525 config-identical paired cells per
motion model, the aggregate ratio moved by $\times1.15$--$1.41$ at every $\rho>0$
and by $\times0.96$--$0.99$ frozen, which is what a defect reachable only through
re-preparation predicts: a frozen lane never re-prepares. Under uniform jump at
$\rho{=}0.25$ the number of operators on which the block-scheduled arm is the
better one went from 9 of 21 to 12 of 21, so an aggregate that read as a minority
result became a majority one. On the 84-cell single-seed grid an earlier version
of this study used, nine hindsight labels changed: eight toward the
block-scheduled arm and one the other way. The defect therefore moved the
classes a predictor is fitted to separate, not only the noise around them.

\paragraph{Change detection, charged to the other family.}
The second one runs the other way, which is why it is worth printing beside the
first: finding one defect that flatters a system is not evidence that the
remaining ones do too.

Both families must detect that their own precondition moved --- the column-exact
arm asks whether $D$ changed at all, the block-scheduled arm asks whether the
dirty \emph{tile} set changed --- and both questions are $O(|D|)$. In the harness
they were not. The column-exact side compared column arrays with a memory
comparison, $0.0018$~ms at $|D|{=}256$; the block-scheduled side built its tile
key as a set comprehension in the harness language, $0.0519$~ms, a $29\times$
asymmetry on cells whose entire kernel is $0.02$--$0.09$~ms. Because the dirty
columns are already sorted, deduplicating their tile indices needs one neighbour
comparison and no set and no sort: $0.0043$~ms, within $2.4\times$ of what the
baseline pays for the same kind of question.

Removing it moves the block-scheduled arm's standing substantially and in its
favour. Over 45 paired cells --- five operators spanning the density range,
three churn rates, three seeds, identical in everything but that one flag --- the
block-scheduled margin over the column-exact baseline rises by $1.48\times$ for
the open reference arm and $1.47\times$ for the proprietary one. The correction
is present frozen as well as churning ($1.55\times$ at $\rho{=}0$, $1.52\times$
at $\rho{=}0.01$, $1.39\times$ at $\rho{=}0.25$), which is what a per-step cost
that does not depend on whether anything moved should do, and it is large enough
to change signs: on \texttt{bcspwr10} the open arm reads $0.86\times$ with the
harness's own key and $1.35\times$ without it, a loss and a win for the same
kernel on the same workload.

Every churn number in this paper that predates the fix is conservative by
roughly that margin, and the direction is the opposite of the
session-construction defect's. That is the reason both are reported here rather
than only the first: finding one defect that happens to flatter the system under
test is not evidence that the remaining ones do.

Stated as a rule rather than as two anecdotes: time what the \emph{caller} pays,
and for each thing timed, check whether the other arm pays it too and in a
comparable implementation. Both defects here passed the first test and failed
the second, and a benchmark that only applies the first will report them as
results.

\section{Motion locality against index locality}
\label{sec:control}

Section~\ref{sec:families} varies one label and reports a large effect, and the
label confounds two things. Local drift refills from the \emph{structural}
neighbourhood of the survivors, while the block-scheduled arm re-prepares when a
column leaves the \emph{column-index} blocks it already covers. Those coincide
only because these operators arrive in orderings that happen to preserve
locality, so 418 of 418 could be a restatement of the generator rather than a
result. This section separates the two, and reports a smaller claim than
\S\ref{sec:families} on its own would support.

\paragraph{The control.}
We relabel each operator, $A' = A[p][:,p]$ for a random permutation $p$, and
replay the identical structural trajectory in the new index space. The graph is
unchanged, the motion model is unchanged, and index locality is gone. The
permutation has to be symmetric. Both the neighbourhood generator and the drift
model walk the operator by using a column index as a row index, so under a
column-only permutation that walk would follow edges into a scrambled index
space and ``refills from the structural neighbourhood'' would stop being true of
the permuted run, confounding the loss of index locality with the loss of the
motion model itself. Under a symmetric relabelling the walk is the original walk
with its nodes renamed.

Permutation seeds drawn: \textbf{4} (101, 202, 303, 404). Each is an independent
symmetric relabelling replayed across all 21 operators, all five dirty-set seeds
and the same $\rho$ grid, so the four runs differ in the permutation and in
nothing else, and the churning cell set is the same 418 cells in all of them.
Relabelled numbers in this section are a median over the four with the range
across them, except where a permutation is named; Figure~\ref{fig:motion} plots
seed 101, the paper's original draw.
Four draws of one relabelling process on one operator set are not four
replications of the experiment. They share the operators and the structural
trajectory, so they establish where the point estimate sits and whether any
single draw was unrepresentative, and they do not tighten the intervals: each
permutation is bootstrapped separately and the four intervals are reported side
by side. Pooling the cells would divide the interval width by roughly two, on
replication the design does not have, and it is interval width rather than the
point estimate that the unresolved contrast below turns on.
A random relabelling is also an extreme rather than a distribution --- real operators are neither
optimally ordered nor uniformly scrambled --- so the control brackets the effect
without saying where a given deployment's operator sits inside the bracket
(\S\ref{sec:limits}).

\paragraph{Result.}
The block-scheduled arm does not still win 418 of 418 under permuted drift. It
wins a median 382 of 418 across the four permutations (range 373--386) at a
geomean of $2.35\times$ (range $2.34$--$2.37$) against $5.62\times$ native
(Figure~\ref{fig:motion}, right). The original draw is the top of both ranges,
386 cells and $2.37\times$, so the number the paper first reported was the most
favourable of the four. Index-tile locality is doing real work and
\S\ref{sec:families} should not be read as a claim about motion alone.

That leaves three comparisons and they do not all carry the same weight. Each is
a paired bootstrap over the 21 operators, 4000 resamples, 95\% CI, run
separately within each permutation, on the unit and with the aggregation rule
the selection contrasts of \S\ref{sec:contrast} use; margins are compared as a
ratio of geomeans and win rates as a difference of proportions. Each is reported
as the median point estimate over the four permutations, the widest interval any
of the four produced, and how many of the four exclude the null.

\begin{itemize}
\itemsep1pt
\item Destroying \emph{index} locality costs the margin $2.39\times$
$[1.93, 2.93]$ ($5.62 \to 2.34$--$2.37$) and 32--45 of 418 cells, $+0.086$
$[+0.000, +0.210]$ on the win rate. The margin effect is large and excludes the
null in 4 of 4 permutations. The cell-count effect excludes it in 3 of 4, and
the exception is seed 101, whose lower bound is $+0.000$: on that metric the
original draw was the only one of the four that reached the null.
\item Destroying \emph{motion} locality costs $2.92\times$ $[2.43, 3.49]$
($5.62 \to 1.92$) and 105 cells, $+0.253$ $[+0.122, +0.399]$. Both exclude the
null. Neither lane in this contrast is relabelled, so it is identical in every
permutation.
\item Local drift with \emph{no} index locality against uniform jump with index
locality intact is the comparison the two are then ranked by, and it is the weak
one. On the win rate it holds in all four: $89.2$--$92.3\%$ against $74.7\%$,
$+0.167$ $[+0.050, +0.281]$, every interval excluding $0$. On the margin it
holds in none: $1.22\times$ $[0.96, 1.57]$, every interval covering $1$. At the
over-corrected end below, neither metric survives in any permutation:
$1.05\times$ $[0.80, 1.37]$ on the margin and $+0.038$ $[-0.12, +0.15]$ on the
win rate, where one permutation's point estimate is itself negative
($-0.015$).
\end{itemize}

The residual is bounded from the wrong side. Relabelling slows the column-exact
arm as well, by a per-permutation median of $1.152$--$1.173\times$ (per cell,
$0.822$--$3.033\times$ over all four), against $2.67$--$2.75\times$ for the
block-scheduled arm, so those cell counts are measured against a baseline that
also degraded. Charging the block-scheduled arm its relabelling penalty while
holding the baseline at its \emph{native} time over-corrects, since it refunds a
cost the baseline really paid, and gives a median 328 of 418 (range 306--333) at
$2.03\times$ (range $1.97$--$2.05$). The isolated effect sits between the two,
so the control's result over four permutations is the bracket
$1.97$--$2.37\times$ and 306--386 of 418 rather than either end. In the original
draw every relabelled loss falls on three operators
(\texttt{scircuit} 12, \texttt{wang3} 11, \texttt{rajat03} 9); those three carry
losses in all four permutations, joined by \texttt{bcspwr10} in three and
\texttt{epb2} in one, for a union of five. All five already
change sides between drift and jump in the main sweep, so in every permutation
relabelling costs the block-scheduled arm a subset of the operators uniform
motion costs it rather than a new class.

So what the control establishes is one thing and not two. Destroying index
locality costs 32--45 of 418 cells and takes the margin from $5.62\times$ to
$2.34$--$2.37\times$: a large effect, it holds in every permutation, and
\S\ref{sec:families}'s 418 of 418 is partly a property of the orderings these
operators ship in. Whether motion locality contributes anything \emph{beyond}
index locality is still not resolved. The win-rate half of that ranking holds in
all four permutations as measured and fails in all four at the over-corrected
end; the margin half fails in all four at both. Replication moved that contrast's
point estimate by $0.015$ and left all four intervals covering $1$, which is what
it should do: more permutations of one operator set narrow nothing, and the
ranking needs a wider operator set. What
does not depend on the ranking is that ``there is no arm-selection problem at
all under local drift'' is true of these operators as they ship rather than of
local drift as such: relabelled, a small selection problem appears, and the
online policy's agreement with hindsight falls from $100\%$ to
$86.8$--$90.7\%$.

\begin{figure}[tbp]
  \centering
  \includegraphics[width=\linewidth]{fig_motion_control}
  \caption{Left: the mixing axis. A fraction $\alpha$ of each step's refills is
  drawn uniformly, so local drift and uniform jump are $\alpha{=}0$ and
  $\alpha{=}1$. Both the margin and the share of cells the block-scheduled arm
  wins fall monotonically, and most of the fall has happened by
  $\alpha{=}0.25$. 21 operators, three seeds, 200 steps; read for shape rather
  than for level (\S\ref{sec:limits}). Right: the four lanes the control
  produces, churning cells only, labelled with the geomean over operators and
  the count of cells the block-scheduled arm wins. The relabelled bars are
  permutation seed 101 of the four \S\ref{sec:control} reports. The gap between
  the first bar and the second is what destroying index locality costs,
  $2.37\times$ $[1.93, 2.90]$ in this permutation and a median $2.39\times$ over
  the four. The gap between the second or third bar and the fourth is the
  contrast \S\ref{sec:control} declines to draw a conclusion from: its interval
  covers $1$ in every permutation.}
  \label{fig:motion}
\end{figure}

\paragraph{The mixing axis.}
\label{sec:mixing}
$\alpha$ makes the contrast a curve rather than a dichotomy
(Figure~\ref{fig:motion}, left). Re-run at the main sweep's own length and seed
count --- 21 operators, five seeds, 400 steps, $\rho \in \{0.01, 0.25\}$ --- the
geomean falls $4.95 \to 3.43 \to 2.87 \to 2.74 \to 1.71$ across
$\alpha \in \{0, 0.25, 0.5, 0.75, 1\}$, and the cells the block-scheduled arm
wins fall $199/209 \to 188 \to 166 \to 169 \to 143$. The number of operators it
wins on \emph{every} cell falls $19 \to 18 \to 15 \to 15 \to 12$ of 21.

\textbf{The shape claim an earlier version of this section made does not
survive that re-run, and did not really survive the earlier one either.} The
claim was that the curve is concave and that most of the distance to uniform
jump is covered by $\alpha{=}0.25$. Measured as a fraction of the total fall in
log space --- the right frame for a ratio of geomeans --- $\alpha{=}0.25$ covers
$34.5\%$ of it here and $39.1\%$ in the 200-step file. A third is not most, in
either.

What the curve actually does is fall at both ends and sit almost flat in the
middle: the linear steps are $-1.52, -0.56, -0.13, -1.03$. Part of that plateau
is an artifact the axis cannot avoid, and \S\ref{sec:motion} names it: $\alpha$
scales an integer count of replacements, so at $\rho{=}0.01$, where
$\lceil \rho|D| \rceil$ is 3, $\alpha{=}0.5$ and $\alpha{=}0.75$ both retire two
columns uniformly and are the same workload measured twice. Those 105 cell pairs
disagree by a geomean of $0.926$ on the block-scheduled lane, which is a floor
measurement rather than a curve feature. The honest statement is the weaker one:
a quarter of each step's replacements drawn uniformly already costs about a
third of the way to uniform jump, so ``mostly local'' motion is meaningfully
worse than local --- but the axis does not collapse onto its jump endpoint
early, and no claim in this paper should rest on its middle.

Split by rate, the selection problem appears first at high churn: at
$\rho{=}0.25$ the block-scheduled arm wins 99 of 105 cells at $\alpha{=}0$ and 90
at $\alpha{=}0.25$, where at $\rho{=}0.01$ the same step costs it 2 (100 to 98).

Even at 400 steps this sweep's levels are not interchangeable with
\S\ref{sec:families}'s, for a different reason than before: it covers two churn
rates rather than four and scores the router's oracle rather than a direct
comparison, so its $\alpha{=}0$ column reads 199 of 209 where \S\ref{sec:families}
reads 418 of 418. The axis is for shape; \S\ref{sec:families} is for level.

\section{What a merged schedule is worth}
\label{sec:merge}

The motion model also decides the value of a schedule optimization that has
nothing to do with arm selection. A delta schedule can be emitted per
column-tile, or it can be consolidated so that each distinct output row appears
once and accumulates contributions from every dirty tile that touches it. The
consolidated form touches fewer distinct output rows, and how many fewer is the
merge's structural value. We report it as $\mathrm{rows}_{\text{per-tile}} /
\mathrm{rows}_{\text{merged}}$, an integer ratio read off the built schedule
rather than a timing, so the measurement floor does not enter it. It is
invariant in batch width, checked on all 15 (operator, seed) cells that carry
three batch widths.

The sweep is 21 operators $\times$ 5 churn rates $\times$ 5 seeds $\times$ both
motion models, 200 steps, $B{=}8$, and the reported cell is the last step. Its
seeds are 17--21, the same five the main sweep uses. An earlier version of this
study covered five operators at seeds 17/117/217; replaying that configuration
reproduces its 120 committed rows exactly, both integers in each, so the change
in level below is the operator set and nothing else.

\paragraph{The row saving is now reproducible, and so is what it is worth.}
The row counts in this section come from schedules only the proprietary arm
builds, which made it the one part of the paper a reader could not recompute. The open
reference arm of \S\ref{sec:families} emits both forms, so the ratio is now an
integer identity anyone can check --- and it agrees across the whole table, not
only at one point. Replaying the trajectory to the last step, exactly as this
section's own convention requires, the open schedules save $2.97\times$ row
touches frozen, $2.76\times$ under local drift and $1.72\times$ under uniform
jump at $\rho \geq 0.01$, against $2.97$, $2.87$ and $1.74\times$ for the
proprietary ones. The frozen figure is identical, the uniform-jump figure agrees
to $1\%$, and the local-drift figure to $4\%$. The separation between the two
motion models --- the thing this section is actually about --- is $1.61\times$
on open code against $1.65\times$ on closed.

Emitting both forms also prices something the row count cannot. A merged
schedule must be re-emitted whenever the live tile set moves, where a per-tile
one compiles only the tiles it has never seen, and the open arm lets both be
timed on identical cells. Frozen, the merge is worth $1.15\times$ --- a real
gain, but far less than the $2.97\times$ row saving suggests, because a row
touch is not the unit that costs. Once the set moves it inverts hard: the merged
form runs at $0.28\times$ the per-tile form at $\rho{=}0.002$ under local drift,
falling to $0.07\times$ at $\rho{=}0.25$, and sits at $0.05\times$ throughout
under uniform jump.

So the merge's structural value and its realized value point in opposite
directions across exactly the boundary this paper is about, and the gap is an
order of magnitude wide. A merged schedule is worth having only if it can be
updated without being rebuilt --- which is the design Appendix~\ref{sec:patch}
builds and measures, and this is the open evidence for why that was the
interesting thing to try. Read on its own, the $2.97\times$ row saving is the
kind of number a benchmark reports as an optimization's value and a deployment
never sees.

\paragraph{Level and separation.}
The merge is worth $2.97\times$ on a frozen set. On a moving one it is
$2.87\times$ under local drift against $1.74\times$ under uniform jump, both
restricted to $\rho \geq 0.01$ so that the two sides cover the same churn rates
(Figure~\ref{fig:merge}, left; 21 operators, five seeds). Taking local drift over
every churning rate instead, which is what a reader of the abstract of an earlier
version of this study saw, reads $2.90\times$ and a separation of $1.66\times$
rather than $1.65\times$; the choice does not matter here and it is stated
because it does matter for the per-operator table below. Restricted to the
original five operators the same file reads $2.16$/$2.13$/$1.26$, which is the
five-operator study it replaces. The separation ratio is stable at
$1.65$--$1.70\times$ across both populations; the levels are not, and the
five-operator numbers understate the merge on a representative operator set by
$35$--$38\%$.

\begin{figure}[tbp]
  \centering
  \includegraphics[width=\linewidth]{fig_merge_value}
  \caption{What consolidating a schedule by global row is worth. A merge pays
  only for output rows that more than one dirty tile contributes to, so this
  ratio measures how much row sharing the motion leaves behind. Left: geomean
  over 21 operators, median over five seeds, both motion models, never averaged.
  Right: the same operators one by one, local drift against uniform jump at
  $\rho \geq 0.01$. Points on the diagonal are operators the motion model does
  not move; points below it are the two on which the merge is worth more under
  uniform jump.}
  \label{fig:merge}
\end{figure}

\paragraph{Uniformity across operators.}
Per-operator sensitivity is drift over jump at matched $\rho \geq 0.01$, under
the aggregation rule of \S\ref{sec:metrics}: median across the five dirty-set
seeds within an (operator, $\rho$) cell, then geometric mean over $\rho$. Of the
21 operators the merge collapses on 13 (sensitivity $\geq 1.5$; \texttt{olafu}
$3.29$, \texttt{raefsky3} $3.00$, \texttt{ct20stif} $2.96$, \texttt{bcsstk25}
$2.32$, and nine more), is moderate on 3, is indifferent on 3
(\texttt{bcspwr09} $1.02$, \texttt{circuit\_3} $1.03$, \texttt{bcspwr10}
$1.05$), and \emph{inverts} on 2: on \texttt{msc10848} ($0.79$) and
\texttt{nasa2910} ($0.85$) the merge is worth more under uniform jump than under
local drift. Both are FEM operators whose merge factor is high in every regime.
These are integer row counts rather than timings, so they are not a measurement
artifact, and nothing in the account below predicts them. The mechanism
paragraph should not be read as covering all 21 operators.

One operator changes band with the convention, and it is worth saying which. A
first version of this table pooled the five seeds flat instead of taking their
median, and took local drift over all churning $\rho$ against uniform jump at
$\rho \geq 0.01$; that reads 14 collapses and 2 moderate, with
\texttt{TSOPF\_RS\_b300\_c1} at $1.55$ rather than $1.28$. The quoted values move
by $2$--$5\%$ ($\texttt{ct20stif}\ 3.09$, $\texttt{raefsky3}\ 3.11$,
$\texttt{scircuit}\ 1.95$ under the unmatched convention). The two
counter-examples are the same two under both.

\paragraph{Mechanism.}
The mechanism shows in the $\rho$ dependence rather than in the model labels. A
merge consolidates rows that more than one dirty tile contributes to, so its
value falls as far as the dirty tiles are scattered, and scattering takes time.
At $\rho{=}0.002$ with $|D|{=}256$ the model retires one column per step, so a
200-step run turns over 200 of 256 positions and the starting neighbourhood is
still a majority of the set through most of it; on 7 of the 21 operators the
uniform-jump merge factor at that rate is still within $5\%$ of its local-drift
value. The
same prediction holds within a single run: measured across the sampled steps,
the merge at step $\geq 150$ is $0.978\times$ its value at step $\leq 25$ under
local drift and $0.750\times$ under uniform jump, so under uniform redraw the
merge is still falling when the run ends and under drift it is flat.

What does \emph{not} hold is any structural predictor of which operators are
sensitive. Over the 21 operators, sensitivity correlates $+0.20$ (Pearson) and
$+0.10$ (Spearman) with tile density. An earlier version of this study read the
flatness of \texttt{bcspwr10} as a low-density effect; that survives as a
statement about those operators --- the two with $\mathrm{dens} < 0.30$ are two
of the three flattest --- and fails as a rule, since the third is
\texttt{circuit\_3}, which sits between them at $\mathrm{dens} = 0.80$, and the
rest of the set has median sensitivity $1.85$. Two better candidates exist, the
number of tiles the dirty set already spans ($-0.47$ Pearson, $-0.46$ Spearman)
and $N$ ($+0.21$ Pearson, $+0.63$ Spearman), but at
21 points with several candidates tried, no multiple-comparison correction and
no held-out set, both are hypotheses.

What is measured here is row sharing and its dependence on the motion model.
The attribution to topological adjacency specifically is inferred: the row
counts do not distinguish adjacency from any other property local drift
preserves and uniform jump destroys. The contrast is controlled --- identical
operators, $|D|$, $\rho$ and seeds --- and it shows in the raw counts, with
\texttt{ct20stif} touching 1{,}892 distinct rows under local drift at
$\rho{=}0.002$ against 11{,}906 under uniform jump.

\section{A cost model, and why one structural predictor had to fail}
\label{sec:model}

Everything above is empirical, which leaves a reader unable to map their own
workload onto it without running the harness, and leaves \S\ref{sec:select}'s
negative result about structural prediction looking broader than it is. A short
cost model fixes both. Per step, with $|D|$ dirty columns at batch width $B$:
\[
  T_{\text{ce}} \;=\; a\,\mathrm{nnz}_e B \;+\; b\, r_e B \;+\; q\,C_e,
  \qquad
  T_{\text{bs}} \;=\; a\,\mathrm{nnz}_s B \;+\; b\, r B \;+\; s\,|D|B \;+\; m\,C_t,
\]
where $\mathrm{nnz}_e = \mathrm{nnz}(A[:,D])$ is what a column-exact arm touches,
$\mathrm{nnz}_s = \mathrm{nnz}(A[:,\mathrm{tiles}(D)])$ is what a tile walk
scans, $\mathrm{dens} = \mathrm{nnz}_e/\mathrm{nnz}_s$, $r_e$ and $r$ are the
distinct output rows each emission touches, $q \in \{0,1\}$ records whether $D$
changed this step, $m$ counts tiles compiled this step, $C_e$ and $C_t$ are the
cost of emitting a whole delta and one tile segment, and $s|D|B$ is the scatter.
Everything but the four machine constants is structural and free to compute.

\paragraph{The motion model enters in exactly one place.}
It is $m$. An arriving column either lands in a tile the arm has already
compiled, which is free, or in one it has not. Writing $\varphi$ for the
probability of the latter,
\[
  \mathbb{E}[m] \;=\; \rho\,|D|\,\varphi ,
\]
and $\varphi$ is a joint property of the motion model and the operator's
topology, measurable from the dirty-set trajectory with no timing at all. Under
local drift an arrival is a structural neighbour of a survivor and neighbours
mostly lie in tiles the survivors already occupy, so $\varphi$ is small; under
uniform jump it approaches the uncovered fraction of live tiles. The paper's
whole headline is that one number.

\paragraph{Prediction 1: the frozen half needs no motion model.}
At $\rho{=}0$ both $q$ and $m$ vanish and the ratio collapses to the apply
alone, so the block-scheduled arm should lose by roughly the padding factor.
Measured on the open reference arm over all 42 frozen cells, the frozen ratio
divided by $\mathrm{dens}$ is $0.59\times$ (range $0.37$--$1.10$): padding is the
right first-order account and explains rather more than half of the gap, and the
operators with $\mathrm{dens}\!\approx\!1$ are the ones where the frozen gap
closes to a tie. What the residual is remains unisolated, here and in
\S\ref{sec:families}.

\paragraph{Prediction 2: \texorpdfstring{$\mathrm{dens}$}{dens} alone had to
fail.} Appendix~\ref{app:selection} reports that $\mathrm{dens}$ --- the
mechanistically motivated structural feature --- separates neither motion
model's classes, and reports it as a negative result about structural
prediction. The model says something narrower and more useful:
$\mathrm{dens}$ prices the apply and says \emph{nothing whatever} about
re-preparation, which is the other half of the cost and the only half the motion
model touches. A feature that is constant in $\rho$ and identical under both
motion models cannot separate an outcome that is neither. The failure was
predictable from the cost structure rather than evidence that no structural
predictor exists, and the model names the missing feature: $\varphi$, which is
obtained the same way $\mathrm{dens}$ is and at comparable cost.

\section{Selecting the arm}
\label{sec:select}

Where neither arm dominates a deployable system must choose, and
\S\ref{sec:families} locates that scope: uniform jump. This section asks what
such a system has to observe in order to choose well. We compare offline
models fitted to features available before running anything against an online
policy that measures both arms on the caller's own workload.

Appendix~\ref{app:selection} carries the protocol, the model table and the
contrasts; this section states where a selection problem exists and what the
comparison returns.

\paragraph{Where a selection problem exists.}
Restricted to churning cells under local drift the labels collapse: the
block-scheduled arm is correct in all 418, so every method --- constant arm,
learned model, online policy --- ties at zero regret by construction, and the
study reports the collapse rather than a contrast. On the workload mixture
(frozen and churning cells scored together, and named as a mixture wherever it
appears) the decision is real but trivial: a depth-1 tree on the observed $\rho$
alone reaches $0.994$ leave-one-operator-out accuracy and $1.000$ geomean
regret, worst cell $1.01\times$, splitting at $\rho > 0$
(Table~\ref{tab:predictors}). The best structural
model reaches $0.805$ and $1.159$, worst cell $5.85\times$, which is exactly the
score of always choosing the block-scheduled arm. No structural model beat that
constant. Under local drift the deployment rule needs no model: switch when
the set moves.

The scope in which the question is live is uniform jump. There, on churning
cells, $\rho$ carries no information: a depth-1 tree on $\rho$ alone scores
$0.747$, identical to always choosing the block-scheduled arm, while structure-only
logistic regression reaches $0.833$ and the best structure-plus-$\rho$ model
$0.888$. So $\rho$ decides \emph{whether to switch arms at all} and does not
decide \emph{which} arm once churn is established. A single-motion-model, single-seed
sweep without a $\rho{=}0.002$ row cannot distinguish those two statements, and
an earlier version of this study did not.

\begin{table}[tbp]
  \centering
  \caption{Speedup is the geomean over always-column-exact; ``worst cell'' is the
  largest slowdown against the per-cell hindsight oracle. Four blocks, never
  collapsed. Under local drift, always-block-scheduled \emph{equals} the oracle
  on churning cells, because it is never the wrong choice there.}
  \label{tab:headline}
  \small
  \begin{tabular}{lcccc}
\toprule
\multicolumn{5}{l}{\emph{local drift}}\\
\cmidrule(lr){1-5}
& \multicolumn{2}{c}{frozen ($n{=}105$)} & \multicolumn{2}{c}{churning ($n{=}418$)}\\
\cmidrule(lr){2-3}\cmidrule(lr){4-5}
strategy & speedup & worst cell & speedup & worst cell \\
\midrule
always column-exact & 1.00$\times$ & 1.01$\times$ & 1.00$\times$ & 35.57$\times$ \\
always block-scheduled & 0.48$\times$ & 5.85$\times$ & 5.57$\times$ & 1.00$\times$ \\
measured router (steady state) & 0.96$\times$ & 1.21$\times$ & 5.12$\times$ & 3.09$\times$ \\
\textbf{measured router (incl. calibration)} & 0.90$\times$ & 1.26$\times$ & 4.00$\times$ & 3.68$\times$ \\
\midrule
hindsight oracle \emph{(not achievable)} & 1.00$\times$ & 1.00$\times$ & 5.57$\times$ & 1.00$\times$ \\
\midrule
\multicolumn{5}{l}{\emph{uniform jump}}\\
\cmidrule(lr){1-5}
& \multicolumn{2}{c}{frozen ($n{=}105$)} & \multicolumn{2}{c}{churning ($n{=}419$)}\\
\cmidrule(lr){2-3}\cmidrule(lr){4-5}
strategy & speedup & worst cell & speedup & worst cell \\
\midrule
always column-exact & 1.00$\times$ & 1.03$\times$ & 1.00$\times$ & 19.92$\times$ \\
always block-scheduled & 0.47$\times$ & 5.87$\times$ & 1.92$\times$ & 2.40$\times$ \\
measured router (steady state) & 0.96$\times$ & 1.21$\times$ & 1.90$\times$ & 1.91$\times$ \\
\textbf{measured router (incl. calibration)} & 0.90$\times$ & 1.27$\times$ & 1.55$\times$ & 2.75$\times$ \\
\midrule
hindsight oracle \emph{(not achievable)} & 1.00$\times$ & 1.00$\times$ & 2.06$\times$ & 1.00$\times$ \\
\bottomrule
\end{tabular}

\end{table}

Table~\ref{tab:headline} gives the results. On churning cells the policy reaches
$5.12\times$ over always-column-exact in steady state under local drift and
$4.00\times$ with all calibration charged, against a $5.57\times$ hindsight
ceiling; under uniform jump, $1.90\times$ and $1.55\times$ against $2.06\times$.
As a fraction of the ceiling that is $91.9\%$ and $92.2\%$ in steady state but
$71.8\%$ and $75.3\%$ once calibration is billed. Those last two are properties
of a 400-step harness rather than of the policy, and Appendix~\ref{sec:amortize} gives
the curve. The uniform-jump figure sits on it directly, at $0.753$ for $S{=}400$;
the local-drift figure does not, because the curve's local-drift scope is the
workload mixture rather than churning cells alone. On frozen cells the ceiling is
$1.00\times$, so there is nothing to win, and the policy delivers $0.96\times$
steady and $0.90\times$ calibrated --- the cost of confirming that --- while
spending zero probes in all 210 frozen cells.

\paragraph{Offline selection against online probing.}
Against the online policy, a nested offline model-selection procedure given
structural features and the observed churn rate improves geomean regret by
$0.044$--$0.078$ in all three scopes that carry a live decision, and no single
model named in advance reproduces that. The probe is a one-time cost and
amortizes in session length. Appendix~\ref{app:selection} gives the contrasts
and their intervals and Appendix~\ref{sec:amortize} gives the amortization
curve.

\section{Limitations}
\label{sec:limits}

\paragraph{Availability of the proprietary arm.} One of the two
block-scheduled implementations is proprietary and is released as measured
behaviour rather than source, so a reader cannot re-derive its timings. What
this costs is now much smaller than it was, because it is no longer the only
member of its family in the study: the open reference arm of
\S\ref{sec:families} carries the family claim, reproduces the dichotomy
independently, and makes \S\ref{sec:merge} --- previously not reproducible from
the artifact at all --- an integer identity a reader can recompute. Every
headline in \S\ref{sec:families} and \S\ref{sec:merge} is therefore stated for
an arm that ships with the benchmark, with the proprietary arm reported beside
it as a second data point.

What remains closed is a comparison rather than a claim: how much of the
shipped arm's margin is its engineering rather than its strategy. A reader can
check that the strategy reproduces; they cannot check our number for that
particular product. We consider this the right residue --- the paper's
contribution is the task, the motion model and the protocol, none of which
depend on it.

\paragraph{Permutation seeds.} The index-locality control of \S\ref{sec:control}
replays one structural trajectory in two index spaces, over four permutation
seeds. Four draws of one relabelling process on one operator set show that no
single permutation was unrepresentative, and they do not narrow the intervals:
each is bootstrapped separately, and pooling the cells would claim an
independence the design does not have. That is why \S\ref{sec:control} reports
the index-locality effect and still declines to rank the two localities. The
ranking is carried by a $1.02$--$1.23\times$ margin contrast whose interval
covers $1$ in all four permutations, where the index-locality effect itself
excludes $1$ in all four at a median $2.39\times$ $[1.93, 2.93]$. More
permutations of these 21 operators would leave both where they are; a wider
operator set is what the ranking is short of.

\paragraph{Measurement conditions and within-cell replication.} All timings come
from a single Intel i7-1165G7 with turbo disabled and the performance governor
pinned, g++ 13.3, PyTorch 2.13.0+cpu, f64. Ratios should be read as that
machine's ratios. No cell is repeated (Appendix~\ref{sec:timing}), so the paper has two
same-workload floor measurements rather than a within-cell variance, and they
agree on the order of magnitude. The two $\rho{=}0$ blocks disagree by a median
$1.041$ per cell with 18 of 105 over $1.10$ and a maximum of $1.553$. The mixing
sweep supplies a second, independent one: $\rho{=}0.01$ retires three columns per
step, so $\alpha{=}0.5$ and $\alpha{=}0.75$ both retire two uniformly and are the
same workload under two labels; those 63 pairs repeat to a geomean of $0.974$ on
the block-scheduled lane and $0.979$ on the column-exact lane, but $0.61$--$1.46$
and $0.70$--$1.31$ per cell.

That floor sits underneath several effects reported here, and it is not small.
Its median is $1.041$ and its p90 $1.117$, so a contrast of a few percent is one
measurement and not a difference. Most of what this paper rests on is either
much larger (the $2$--$7\times$ margins of \S\ref{sec:families}, and the
$2.37\times$ $[1.93, 2.90]$ index-locality effect of \S\ref{sec:control}) or is
not a timing at all (the integer row counts of \S\ref{sec:merge} and
Appendix~\ref{sec:patch}). Two things are not. The $2.4\%$ arena cost,
single-operator ratios quoted to two decimals and the per-cell win counts at
$\rho{=}0.01$ in the mixing sweep are inside the floor outright. And the
contrast \S\ref{sec:control} uses to rank the two localities against each other
is $1.02$--$1.23\times$, which spans the floor's own width; that is a second reason,
independent of the bootstrap, why \S\ref{sec:control} reports the effect of
destroying index locality and does not rank motion locality above it.

\paragraph{Compiler and machine.} The multi-seed brackets
quantify \emph{dirty-set variance only}. They are not machine brackets and not
toolchain brackets, and those are larger: we measure compiler codegen alone as a
$3$--$7\times$ effect on this kernel. It does not threaten every result here
equally. A $3$--$7\times$ codegen effect
cannot flip the $6.94\times$ family margin or the integer row counts, so the
dichotomy of \S\ref{sec:families} and the merge result of \S\ref{sec:merge} are
toolchain-robust. It can easily flip the 53 uniform-jump churning cells that sit
within $10\%$ of parity, and those cells are where the selection study of
\S\ref{sec:select} does its work, so the selection result is toolchain-fragile
and should be given less confidence than the dichotomy. \S\ref{sec:control}
splits: its index-locality effect is $2.37\times$ and sits with the dichotomy,
while the $1.02$--$1.23\times$ contrast between the two localities sits with the
selection result.

We now measure a second toolchain rather than only conceding the point. The
same source, the same flags and the same 378 cells, built with clang 18 instead
of g++ 13.3:

\begin{itemize}\itemsep2pt
  \item \textbf{The dichotomy is toolchain-robust, as predicted.} The open
    reference arm wins \textbf{126 of 126} churning local-drift cells under
    \emph{both} compilers, and loses on frozen cells under both (9 of 126 won
    under g++, 0 of 126 under clang). The proprietary arm reads 119 and 118 of
    126 on the same cells.
  \item \textbf{The levels are not.} The open arm's local-drift geomean falls
    from $5.17\times$ to $3.83\times$, a $26\%$ move from nothing but the
    compiler --- which is the concrete form of the $3$--$7\times$ codegen
    caveat, and the reason \S\ref{sec:families} reports win counts as its
    reproducible quantity.
  \item \textbf{The open-versus-proprietary gap under uniform jump survives.}
    The open arm wins 115 of 126 uniform-jump churning cells against the
    proprietary arm's 66 under g++, and 97 against 63 under clang. The claim in
    \S\ref{sec:families} that a substantial part of the uniform-jump selection
    problem belongs to one implementation is therefore not an artifact of one
    compiler.
  \item \textbf{One ratio does reverse.} Comparing the two block-scheduled arms
    directly under local drift, $p3/\text{open}$ moves from $0.822\times$ to
    $1.034\times$ --- from the open arm being $22\%$ faster to a tie. The cause
    is isolated and is in one lane: clang costs the open reference arm
    $1.262\times$ while leaving the column-exact baseline ($0.968\times$) and
    the proprietary arm ($1.006\times$) alone. That is a codegen difference in
    one implementation, not a property of either strategy, and it is why
    \S\ref{sec:families} states the drift ranking with this caveat attached.
  \item The structural quantities do not move at all: the merge row-count ratio
    is identical on 378 of 378 cells, as an integer identity must be, and
    maximum relative error under clang is $2.0\times10^{-15}$.
\end{itemize}

What we still do not have is a second \emph{machine class}: every number here
comes from one 4-core 28\,W mobile part, and a server-class memory hierarchy is
the obvious place for a memory-bound sparse kernel's ratios to move. Nothing in
this paper should be read as a claim about that machine.

\paragraph{Scope of the mixing sweep.} $\alpha$ is now swept at the main
sweep's own 400 steps and five seeds, so step count is no longer the reason its
levels differ. Two reasons remain: it covers two churn rates rather than four,
and it scores the router's per-cell oracle rather than a direct arm comparison.
Its $\alpha{=}0$ column therefore reads 199 of 209 where \S\ref{sec:families}
reads 418 of 418, and only its shape is interchangeable. The earlier 200-step,
three-seed version of this sweep is the one an earlier draft's concavity claim
rested on; re-running it did not rescue that claim (\S\ref{sec:mixing}).

\paragraph{A between-block offset in the noise floor.} Across the two frozen
blocks the 21-operator geomean spans $0.473$--$0.479$, and one block sits below
the other on all five seeds, which is a systematic between-block offset rather
than independent draws. Its sign is not stable across harness vintages. It
cannot be a seed effect, since the same seeds are replayed in every block. We
did not establish its mechanism.

\paragraph{A confound in one motion model, and its control.} Uniform jump
scatters the dirty set and, on 8 of 21 operators, also shrinks its nonzero count,
because a random column slice is not the same size as a structural neighbourhood
--- high-degree columns cluster, so a neighbourhood over-samples them and a
uniform draw does not. On the other 13 the slice size barely changes and the
model measures locality alone. Every drift-versus-jump contrast on those 8
therefore compares two changes at once.

The benchmark now ships a third motion model that separates them.
\texttt{jump\_nnz} replaces each retired column with one drawn uniformly from the
live columns whose nonzero count matches its own to within $10\%$, so locality is
destroyed exactly as \texttt{jump} destroys it while the slice's nonzero count is
preserved column by column rather than merely in aggregate --- which matters,
because two slices can share a total and still give the arms different work
depending on how the degree is distributed across tiles.

\textbf{The confound is larger than the sentence above suggests, and it runs
against uniform jump.} Measured as the slice's nonzero count at the last step
over its count at the first, over 400 steps at $\rho \geq 0.01$: local drift ends
at $0.818\times$ its starting work and uniform jump at $0.415\times$, while
\texttt{jump\_nnz} holds $1.000\times$ (range $0.992$--$1.007$ over all 189
cells, against $0.010$--$1.262$ for plain jump). The eight operators named above are the ones this shows up on, and on the worst
of them the effect is not a $20\%$ correction at all:
the worst single cell, \texttt{TSOPF\_RS\_b39\_c30} at $\rho{=}0.25$, ends at
$0.010\times$ its starting nonzero count, so by the end of that run the
uniform-jump lane is doing $100\times$ less work than the drift lane it is being
compared against.

That helps the column-exact arm disproportionately, because its cost scales with
the slice's nonzeros while a tile walk has a floor set by the tiles it must
scan. Holding the work and destroying only locality, over 189 paired cells at
matched operators, churn rates and seeds, the block-scheduled margin under
uniform jump rises from $1.82\times$ to $2.89\times$ for the open reference arm
and from $1.19\times$ to $1.98\times$ for the proprietary one --- a $1.58$ and
$1.66\times$ correction, in the direction that makes uniform jump \emph{less}
damaging than this paper reports.

The correction grows with the churn rate --- $1.40\times$ at $\rho{=}0.01$,
$1.63\times$ at $\rho{=}0.05$, $1.74\times$ at $\rho{=}0.25$ --- which is what a
shrinkage that accumulates over steps predicts, and it is concentrated on the
operators where the shrinkage is: \texttt{memplus}, \texttt{case9} and
\texttt{circuit\_3} each pay a factor of about $6$ in slice nonzeros and see
their margins move by $2.9$--$3.9\times$, while the operators whose slice
barely moves see corrections within $\pm 10\%$.

So the honest reading of every drift-versus-jump contrast here is that it is an
upper bound on what destroying locality costs, and part of what it attributes to
locality is the workload getting smaller. The direction matters for the paper's
own conclusions: it does not touch local drift, and it narrows rather than widens
the gap between the two motion models. We report the uncorrected numbers
throughout, because they are what the released sweep contains and because the
correction is uniform in sign, and we release \texttt{jump\_nnz} so that anyone
can measure the corrected version on their own operators.

\paragraph{Batch width in the selection sweep.} The selection study fixes $B{=}8$.
The family dichotomy is measured at $B\in\{1,8,16\}$ (Figure~\ref{fig:frozen})
and the ordering does not change, but the selection results are single-$B$.

\paragraph{Regime change.} Both motion models change regime at most once. A
workload that alternates faster than a probe can settle would pay the probe tax
repeatedly, which is the scope in which online measurement is most likely to be
worth its cost and the scope this benchmark does not test. The amortization
curve of Appendix~\ref{sec:amortize} extrapolates one regime for the same reason.

\paragraph{Scope of the prediction result.} We do not claim that no predictor can
exist, nor that offline prediction is universally preferable. We claim that under
leave-one-operator-out evaluation on 21 operators, a nested model-selection
procedure given structural features and the observed churn rate beats a
both-arms probe on geomean regret in all three scopes with a hindsight ceiling
above $1.00\times$, that the same holds with the one paid feature deleted, and
that no single pre-registered model reproduces it. On the two frozen scopes the
same contrast covers zero.

\section{Artifact}
\label{sec:artifact}

We release the operator manifest and fetch script, the dirty-set generator (BFS
growth, all three motion models and the mixing parameter), the open column-exact
CSR delta baseline, \textbf{the open block-scheduled reference arm of
\S\ref{sec:families} in both its per-tile and merged forms}, the solver
instrumentation of \S\ref{sec:realmotion}, the benchmark harnesses, the metric
protocol implementation, the complete result CSVs for every sweep in this paper,
the leave-one-operator-out study and its bootstrap, the bootstrap for the
index-locality control of \S\ref{sec:control}, and the figure and table
generators.

Releasing a member of \emph{each} family is what makes the central result
runnable end to end: a reader can reproduce the frozen loss, the churning win,
the merge identity and the two measurement taxes of \S\ref{sec:defect} without
any proprietary component. The reference arm ships with a correctness suite
that checks it against a direct evaluation, checks that a preparation stays
exact for any dirty set inside the tiles it was prepared for, and checks that
its scratch returns to zero --- the three ways an implementation of this shape
can be fast and wrong. Every number in every
figure, table and paragraph of this paper is printed from one of those CSVs by a
committed script, with one exception. The exception is the $0.82$ against
$0.35$~ms/step of Appendix~\ref{sec:timing}: it comes from a retired harness vintage,
its per-step traces are not among the released files, and it is an illustration
of a measurement rule rather than a result the paper rests on. Everything else
is reachable by running \texttt{paper\_numbers.py}, which tags each value with
the section that quotes it. The offline procedure the paper reports is the one
the released study emits.

We do not release per-step cost traces. Their absence is what makes the selection
results reproducible only as far as the aggregates go, and adding them is the
single change that would most reduce the cost of the closed arm.

\section{Conclusion}

Delta-input SpMM is a widespread workload that existing sparse benchmarks
structurally cannot see, because they hold the operand's relationship to its
predecessor outside the experiment. Letting the dirty set move reveals an
opportunity a one-shot protocol reports as nothing, appearing at a churn of two
columns in a thousand.

Letting it move in two different ways reveals that the opportunity is not one
number. How the dirty set moves decides how large the gap between the two
implementation families is ($6.94\times$ against $2.95\times$ at the same churn
rate), whether an implementation-selection problem exists at all (418 of 418
cells one way, 313 of 419 the other --- though a second block-scheduled
implementation reaches 395 of 419 there, so part of that problem belongs to an
implementation and not to the family), and what a schedule merge is worth
($2.87\times$ against $1.74\times$ over 21 operators at $\rho \geq 0.01$). A
control separates that from the index order the operators happen to ship in and
finds, over four independent permutations, that the index order accounts for
part of it, a median $2.39\times$ $[1.93, 2.93]$ of the $5.62\times$ native
margin; which of the two localities
matters more is not resolved by 21 operators. None of it is visible in the
operator's structure, and none of it is visible in the churn rate alone. A
benchmark for this task has to name its motion model, and a claim about this
task that does not name one is underspecified.

Four additions decide how far any of that generalizes, and each is a different
kind of check. An independently written open implementation of the
block-scheduled strategy reproduces the dichotomy --- a frozen loss, and 418 of
418 churning cells under local drift, the same count the proprietary arm reaches
--- which makes the result a property of the strategy rather than of one
product, and which no amount of care with a single pair of arms could have
established. It also disagrees where it matters: under uniform jump it wins 395
of 419 where the proprietary arm wins 237, so part of what \S\ref{sec:select}
treats as a selection problem belongs to an implementation. Reading the dirty
set off two real solvers rather than generating it finds
$\alpha_{\text{emp}} = 0.000$ for transient propagation on every operator tried
and $\alpha_{\text{emp}} \approx 0.44$ for Newton convergence, so the regime this
paper is about is real for time-stepping workloads and is not the regime an
iterative solve occupies. A motion model that holds the dirty slice's work while
destroying its locality shows that uniform jump was being charged for both, and
that every drift-versus-jump contrast here is an upper bound by $1.58\times$.
And two measurement taxes, pointing in opposite directions, each moved an
aggregate comparison enough to change what it said.

That last pair is the part most likely to matter to somebody else's benchmark.
Both defects passed the obvious test --- time what the caller pays --- and failed
the one that is easy to skip: check whether the other arm pays it too, in a
comparable implementation. A delta benchmark compares implementations that
prepare differently, so it is structurally prone to charging one of them for
work the other never does, and we found two in our own harness after several
rounds of review.

Where a selection problem does exist, the cheapest sufficient observation wins.
An offline model-selection procedure given free structural metadata and the
observed churn rate beats a policy that probes both arms, because the probe has
to run the losing arm and the losing arm loses by a great deal. The probe does
amortize: the policy holds $92$--$93\%$ of the hindsight ceiling on a long
session rather than the $75$--$78\%$ a 400-step harness reports, and it still does not overtake
the nested procedure at any session length. It does overtake the pre-registered
logistic regression, at $S \approx 1000$--$1300$ on the two mixtures
(Appendix~\ref{sec:amortize}).
That leaves online measurement with two arguments this benchmark does
not settle: regime change, which it does not test, and arm sets for which no
cheap scalar exists.

\bibliographystyle{plainnat}
\bibliography{refs}

\appendix

\section{Protocol and validation details}
\label{app:protocol}

\subsection{Timing protocol}
\label{sec:timing}

A cell is one (operator, $\rho$, seed, motion model) at $B{=}8$, run for 400
steps. The dirty-set stream, the $dx$ blocks and the column index tensors are
generated once per cell, outside every timed region, and handed to every
implementation, so all of them replay the same workload step for step. Each step
is timed individually and the timed region covers exactly what a caller would
pay at that step: the apply, plus any re-preparation that fell there. The cell's
reported cost is the \emph{mean} over the 400 steps, not the median, because
re-preparations are rare and lumpy and a median hides them.

Implementations run in a fixed order within a cell (column-exact,
block-scheduled, patchable, then the policy). The three-arm sweep of
\S\ref{sec:families} is the exception and rotates its five lanes by step, which
is one reason its levels are not interchangeable with this one's; the ordering
effect itself was measured on three operators and moved no ratio by more than
$5\%$. Each implementation's \emph{first}
preparation is charged as startup and excluded from the per-step stream: it is
paid once whatever the workload does, and charging it as a churn rebuild priced
the baseline at $0.82$~ms/step against its true $0.35$ on one operator. That
last pair is from a retired harness vintage and is the one number in this paper
not reachable from the released CSVs (\S\ref{sec:artifact}).
No cell is repeated. The replication in this benchmark is the five dirty-set
seeds, which are five different workloads rather than five measurements of one;
\S\ref{sec:limits} reports two separate same-workload floor measurements and
they are the number every ratio here has to be read against.

The one median in the paper belongs to the online policy rather than to the
harness: its arm estimator is
$\widehat{c} = \mathrm{median}(\text{apply}) + \widehat{r}\cdot
\mathrm{mean}(\text{re-prepare})$ over its probe window (\S\ref{sec:measure}).

\subsection{Exactness}

Absolute error is not a usable criterion here. Entry magnitudes span many orders
of magnitude across this operator set, so on the largest-valued operators a
last-bit-correct result reads as a large absolute error, and the baseline reads
the same. Maximum \emph{relative} error over the 105 distinct exactness checks in
the sweep is $1.6\times10^{-15}$, and $8.4\times10^{-16}$ excluding the single
operator that sets it. That figure is an upper bound taken across all lanes
against one common reference, so it does not attribute the largest value to any
particular implementation. Any
implementation that trades accuracy for speed is out of scope; the task is exact
maintenance.

The open reference arm is held to the same standard and checked more finely,
because it is new code and an implementation of its shape can be fast and wrong
in three specific ways. Its released test suite checks the product against a
direct evaluation at $B \in \{1,2,4,8,16\}$ on operators spanning the density
range; that a preparation made for one dirty set stays exact when applied to a
\emph{different} dirty set drawn from the same column tiles, which is the
family's defining property and the thing the churn measurement leans on, and that
no tile is recompiled when that happens; and that its scatter buffer returns to
all-zero after every apply, since a scatter that forgot to un-scatter would leave
the previous step's operand in place and produce a wrong answer only once the
dirty set moved --- exactly the regime being measured. Over the 1050-cell
three-arm sweep the maximum relative error across every lane is
$2.0\times10^{-15}$, the same order as the two-arm bound above.

\subsection{Validation of the index-locality control}

Four checks, recorded per cell rather than asserted, and re-run within each of
the four permutations rather than carried over from the first. The relabelled
operator reproduces the native product to a maximum relative error of
$8.1\times10^{-16}$ to $1.3\times10^{-15}$, depending on the permutation, over
all 525 relabelled cells of each, which fails if the
permutation is applied to the operator but not the dirty set, to the dirty set
but not the operator, or in the wrong direction. $|D|$ and $\mathrm{nnz}$ are
identical on 525 of 525 paired cells in every permutation, so all four index
spaces replay the same structural sets step for step. All implementations agree
within a cell to $1.6\times10^{-15}$, the same bound as the main sweep, in every
permutation. And the native half of
the run reproduces the committed sweep: 418 of 418, per-$\rho$ geomeans
$6.97/6.12/4.91/4.75$ against the main sweep's $6.94/6.15/4.89/4.74$. The native
half does not depend on the permutation and is measured once.

\section{Selecting the arm, in full}
\label{app:selection}

\subsection{The cost of a wrong choice, and a structural predictor}

\paragraph{The penalty for choosing wrongly.}
Mis-routing is expensive and unequally so, and both halves of the asymmetry are
properties of the motion model (Figure~\ref{fig:asym}). On churning cells under
local drift, committing to the column-exact arm costs up to $35.6\times$
($n{=}418$), and committing to the block-scheduled arm costs at most
$1.00\times$ because there is no churning cell where it is wrong ($n{=}0$).
Under uniform jump the two sides are $19.9\times$ ($n{=}313$) and $2.40\times$
($n{=}106$). On frozen cells the block-scheduled side of the penalty reaches
$5.85\times$ (drift, $n{=}102$) and $5.87\times$ (jump, $n{=}100$); those are
frozen numbers, they are not in Figure~\ref{fig:asym}, and they are not
interchangeable with the churning ones.

\paragraph{Structural separability.}
Let $\mathrm{dens}$ be the fraction of scanned entries that are actually needed,
the reciprocal of the padding above, computable from $A$
and $D$ with no timing. It is the mechanistically motivated predictor: it is
what determines the block-scheduled arm's per-step disadvantage. It does not
work in either model, for different reasons (Figure~\ref{fig:overlap}). Under
local drift there is nothing to separate: the block-scheduled arm wins at
$\mathrm{dens}$ from $0.17$ to $1.00$, so any fitted threshold must sit at or
below $0.17$ and the feature carries no information. Under uniform jump the two
classes overlap on $\mathrm{dens} \in [0.17, 0.91]$, an interval containing 340
of 419 churning cells; the best threshold fitted \emph{with hindsight on the test
set} reaches $0.747$ accuracy at $\tau = 0.558$, which is exactly the accuracy of
always choosing the block-scheduled arm.

\begin{figure}[tbp]
  \centering
  \begin{subfigure}{0.48\linewidth}
    \includegraphics[width=\linewidth]{fig_asymmetry}
    \caption{Cost of the wrong choice, churning cells only.}
    \label{fig:asym}
  \end{subfigure}\hfill
  \begin{subfigure}{0.48\linewidth}
    \includegraphics[width=\linewidth]{fig_overlap}
    \caption{The structural predictor against the realized arm.}
    \label{fig:overlap}
  \end{subfigure}
  \caption{Left: mistakes are expensive and expensive unequally, and under local
  drift one direction of mistake does not exist among churning cells. Right:
  under local drift there is no second class; under uniform jump the classes
  overlap on the one feature theory recommends.}
\end{figure}

\subsection{Offline models: protocol}

Eight model families (constant arms, depth-1 and depth-3 decision trees,
logistic regression, $k$-NN, random forest, gradient boosting) under
\textbf{leave-one-operator-out} cross-validation: all 25 rows of an operator
leave together, so no seed of an operator can leak into the fold that scores it.
A random split over cells would place churn rates and seeds of the same operator
on both sides, letting the model memorize the operator, which deployment never
permits. Three feature sets are scored separately, because conflating them
flatters prediction: \emph{structure} (operator-intrinsic only: $\mathrm{dens}$,
$\log N$, $\log \mathrm{nnz}$, mean row nnz, dirty fraction), \emph{observed
$\rho$} (the measured churn rate alone, a count of changed columns rather than a
structural prediction), and \emph{both}. We report accuracy and the distribution
of \textbf{regret}: realized slowdown against the per-cell hindsight oracle.
Every comparison between two strategies is a paired bootstrap over the 21
operators, 4000 resamples, 95\% CI (Table~\ref{tab:contrasts}). Where an offline
strategy is compared against the policy, the offline side is a \textbf{nested}
procedure: the learner is chosen by an inner leave-one-operator-out inside each
training fold, so the choice never sees the rows that score it.

\begin{table}[tbp]
  \centering
  \caption{Offline prediction under leave-one-operator-out cross-validation, per
  motion model, on the workload mixture (frozen and churning cells scored
  together). Accuracy is reported and is not used as an argument.
  \S\ref{sec:contrast} measures what share of an accuracy gap sits on cells
  whose labels this harness cannot resolve; under a flat $10\%$ parity band the
  share is $4.0\%$ on the two local-drift rows and $21.5\%$ on the two
  uniform-jump rows, and it reaches $87$--$100\%$ on uniform-jump churning
  cells, which are a different scope and are not in this table.}
  \label{tab:predictors}
  \small
  \begin{tabular}{llccccc}
\toprule
features & model & LOMO acc. & in-sample & \multicolumn{3}{c}{regret vs oracle}\\
\cmidrule(lr){5-7}
 & & & & geomean & p90 & worst \\
\midrule
\multicolumn{7}{l}{\emph{local drift --- workload mix (frozen + churning)}}\\
\midrule
--- & always column-exact & 0.195 & 0.195 & 3.944 & 22.36 & 35.57 \\
--- & always block-scheduled & 0.805 & 0.805 & 1.159 & 1.87 & 5.85 \\
structure & 1-D threshold (tree d=1) & 0.805 & 0.805 & 1.159 & 1.87 & 5.85 \\
structure & decision tree d=3 & 0.805 & 0.805 & 1.159 & 1.87 & 5.85 \\
structure & logistic regression & 0.805 & 0.805 & 1.159 & 1.87 & 5.85 \\
structure & kNN (k=5) & 0.793 & 0.753 & 1.198 & 2.37 & 13.45 \\
structure & random forest (300) & 0.805 & 0.805 & 1.159 & 1.87 & 5.85 \\
structure & gradient boosting & 0.805 & 0.805 & 1.159 & 1.87 & 5.85 \\
observed $\rho$ & 1-D threshold (tree d=1) & 0.994 & 0.994 & 1.000 & 1.00 & 1.01 \\
observed $\rho$ & decision tree d=3 & 0.994 & 0.994 & 1.000 & 1.00 & 1.01 \\
observed $\rho$ & logistic regression & 0.805 & 0.805 & 1.159 & 1.87 & 5.85 \\
observed $\rho$ & kNN (k=5) & 0.803 & 0.805 & 1.159 & 1.87 & 5.85 \\
observed $\rho$ & random forest (300) & 0.994 & 0.994 & 1.000 & 1.00 & 1.01 \\
observed $\rho$ & gradient boosting & 0.994 & 0.994 & 1.000 & 1.00 & 1.01 \\
both & 1-D threshold (tree d=1) & 0.994 & 0.994 & 1.000 & 1.00 & 1.01 \\
both & decision tree d=3 & 0.985 & 0.996 & 1.003 & 1.00 & 1.40 \\
both & logistic regression & 0.805 & 0.805 & 1.159 & 1.87 & 5.85 \\
both & kNN (k=5) & 0.847 & 0.893 & 1.138 & 1.68 & 5.85 \\
both & random forest (300) & 0.994 & 0.998 & 1.000 & 1.00 & 1.01 \\
both & gradient boosting & 0.985 & 0.998 & 1.003 & 1.00 & 1.40 \\
\cmidrule(lr){1-7}
\multicolumn{2}{l}{\textbf{measured router (steady)}} & 0.994 & --- & \textbf{1.079} & 1.17 & \textbf{3.09} \\
\multicolumn{2}{l}{\textbf{measured router (+calib)}} & 0.994 & --- & \textbf{1.330} & 1.64 & \textbf{3.68} \\
\midrule
\multicolumn{7}{l}{\emph{uniform jump --- workload mix (frozen + churning)}}\\
\midrule
--- & always column-exact & 0.393 & 0.393 & 1.780 & 4.48 & 19.92 \\
--- & always block-scheduled & 0.607 & 0.607 & 1.227 & 2.03 & 5.87 \\
structure & 1-D threshold (tree d=1) & 0.702 & 0.731 & 1.144 & 1.73 & 6.84 \\
structure & decision tree d=3 & 0.683 & 0.769 & 1.152 & 1.74 & 6.84 \\
structure & logistic regression & 0.683 & 0.758 & 1.150 & 1.76 & 4.55 \\
structure & kNN (k=5) & 0.637 & 0.720 & 1.195 & 1.87 & 4.48 \\
structure & random forest (300) & 0.672 & 0.771 & 1.156 & 1.76 & 4.55 \\
structure & gradient boosting & 0.662 & 0.771 & 1.169 & 1.79 & 12.41 \\
observed $\rho$ & 1-D threshold (tree d=1) & 0.788 & 0.788 & 1.056 & 1.24 & 2.40 \\
observed $\rho$ & decision tree d=3 & 0.788 & 0.788 & 1.056 & 1.24 & 2.40 \\
observed $\rho$ & logistic regression & 0.607 & 0.607 & 1.227 & 2.03 & 5.87 \\
observed $\rho$ & kNN (k=5) & 0.578 & 0.424 & 1.250 & 2.24 & 5.87 \\
observed $\rho$ & random forest (300) & 0.788 & 0.788 & 1.056 & 1.24 & 2.40 \\
observed $\rho$ & gradient boosting & 0.788 & 0.788 & 1.056 & 1.24 & 2.40 \\
both & 1-D threshold (tree d=1) & 0.788 & 0.788 & 1.056 & 1.24 & 2.40 \\
both & decision tree d=3 & 0.824 & 0.906 & 1.049 & 1.15 & 3.96 \\
both & logistic regression & 0.679 & 0.759 & 1.152 & 1.76 & 4.55 \\
both & kNN (k=5) & 0.691 & 0.882 & 1.158 & 1.74 & 4.41 \\
both & random forest (300) & 0.880 & 1.000 & 1.030 & 1.05 & 2.03 \\
both & gradient boosting & 0.882 & 0.992 & 1.032 & 1.05 & 3.96 \\
\cmidrule(lr){1-7}
\multicolumn{2}{l}{\textbf{measured router (steady)}} & 0.885 & --- & \textbf{1.076} & 1.19 & \textbf{1.91} \\
\multicolumn{2}{l}{\textbf{measured router (+calib)}} & 0.885 & --- & \textbf{1.281} & 1.64 & \textbf{2.75} \\
\bottomrule
\end{tabular}

\end{table}

\subsection{The online policy}
\label{sec:measure}

The policy treats both arms as black boxes and never inspects the operator. It
opens committed to the column-exact arm, because a dirty set that has not moved
cannot have forced a re-preparation. At the first movement in $D$ it probes:
both arms run on the caller's own workload for a bounded window of 32 steps, and
each arm's steady cost is estimated as
\[
  \widehat{c} \;=\; \mathrm{median}(\mathrm{apply}) \;+\;
                    \widehat{r} \cdot \mathrm{mean}(\text{re-prepare}),
\]
decomposed rather than taken as a window mean because re-preparations are rare
and lumpy: a short window holds one or none, and a raw mean is dominated by
whether the spike happened to land inside it. The cheaper arm is committed to.
Re-probes fire on regime change, and the incumbent keeps serving throughout,
since a re-probe fires when an arm has gone bad and serving with the arm under
suspicion would make every re-probe expensive. Nothing in the policy is specific
to either arm or to sparse linear algebra; it requires only that the arms be
interchangeable and individually timeable.

It picks the hindsight arm in 418 of 418 churning cells under local drift, where
the arms are far enough apart that a 32-step probe separates them every time,
and in 364 of 419 under uniform jump, where the two arms sit close together on
many operators and a probe has less to separate. Frozen agreement is 102 of 105
and 100 of 105. Maximum relative error across the sweep is $1.6 \times 10^{-15}$
including every arm switch; switching arms mid-stream changes timing and nothing
else. Median one probe per run; the maximum is six under local drift and one
under uniform jump.

\subsection{An offline model given the observed churn rate}
\label{sec:contrast}

The comparison that matters is not offline-versus-online, because both sides
measure something. The question is \emph{what} is measured: one scalar, or both
arms. The observed $\rho$ is a count of changed columns and is free --- the
generator retires $\mathrm{round}(\rho|D|)$ columns deterministically, so the
step's count of changed columns recovers $\rho$ exactly, with no estimation
window. The probe must run the losing arm for a bounded window and is not free.
We therefore give an offline procedure the structural features \emph{and} $\rho$,
and score it against the online policy on identical cells.

\paragraph{Procedure, not model.} Five scopes are non-degenerate; three of them
carry a live decision. Against the online policy in steady state the nested
procedure improves geomean regret by $0.078$ $[0.054, 0.101]$ on the local-drift
mixture, $0.060$ $[0.040, 0.081]$ on uniform-jump churning cells, and $0.044$
$[0.019, 0.069]$ on the uniform-jump mixture. All three intervals exclude zero.
On the two frozen-only scopes, where the ceiling is $1.00\times$, the same
contrast covers zero ($-0.027$ $[-0.074, +0.030]$ drift, $-0.003$ $[-0.058,
+0.057]$ jump); those rows are in Table~\ref{tab:contrasts} rather than omitted,
since a claim of the form ``in every scope'' is checkable only if the scopes that
do not support it are printed.

Nesting addresses the objection that the learner is chosen by the same held-out
regret that then scores it. The spread across the six learners is
$0.089$--$0.159$, which bounds the bias, but the bias itself is $0.000$--$0.005$
and the contrast does not move when the selection is nested. What does not
survive is the claim as a statement about \emph{a model}. A logistic regression
pre-registered on principle before any score was read --- lowest-capacity learner
that consumes every feature, monotone decision surface, no hyperparameter to tune
--- \emph{loses} to the policy on the local-drift mixture ($+0.080$ $[+0.047,
+0.114]$) and on the uniform-jump mixture ($+0.076$ $[+0.016, +0.150]$). So the
claim is about a model-selection procedure and not about a model: the procedure
beats the policy in all three scopes, and no single model named in advance does.
A depth-1 tree, equally defensible since the deployed threshold rule is literally
one, does not lose anywhere. ``No single model named in advance'' is the right
statement, and it is narrower than it sounds.

\begin{table}[tbp]
  \centering
  \caption{Paired bootstrap over the 21 operators, 4000 resamples, 95\% CI.
  Positive means the first strategy is worse. Bold marks intervals excluding
  zero. The offline side of every contrast is the nested procedure. Every
  non-degenerate scope is listed, including the frozen ones; local-drift
  churning cells are absent because that scope is single-class.}
  \label{tab:contrasts}
  \small
  \begin{tabular}{lccc}
\toprule
scope & geomean reg. & worst-case reg. & accuracy\\
\midrule
\multicolumn{4}{l}{\emph{local drift}}\\
\midrule
\multicolumn{4}{l}{(a) nested structure-only $-$ policy}\\
\quad frozen ($n{=}105$) & $-0.027$ {\scriptsize$[-0.07,+0.03]$} & $+0.049$ {\scriptsize$[-0.21,+0.21]$} & $-0.048$ {\scriptsize$[-0.14,+0.00]$} \\
\quad mixture ($n{=}523$) & {\boldmath $+0.080$ {\scriptsize$[+0.05,+0.11]$}} & {\boldmath $+2.766$ {\scriptsize$[+0.77,+4.50]$}} & {\boldmath $-0.189$ {\scriptsize$[-0.20,-0.17]$}} \\
\addlinespace[2pt]
\multicolumn{4}{l}{(b) $\rho$-only tree $-$ nested structure-only}\\
\quad frozen ($n{=}105$) & $-0.016$ {\scriptsize$[-0.05,+0.00]$} & $-0.253$ {\scriptsize$[-0.40,+0.00]$} & $+0.048$ {\scriptsize$[+0.00,+0.14]$} \\
\quad mixture ($n{=}523$) & {\boldmath $-0.159$ {\scriptsize$[-0.20,-0.11]$}} & {\boldmath $-4.293$ {\scriptsize$[-4.85,-2.85]$}} & {\boldmath $+0.189$ {\scriptsize$[+0.17,+0.20]$}} \\
\addlinespace[2pt]
\multicolumn{4}{l}{(c) nested structure$+\rho$ $-$ policy}\\
\quad frozen ($n{=}105$) & $-0.027$ {\scriptsize$[-0.07,+0.03]$} & $+0.049$ {\scriptsize$[-0.21,+0.21]$} & $-0.048$ {\scriptsize$[-0.14,+0.00]$} \\
\quad mixture ($n{=}523$) & {\boldmath $-0.078$ {\scriptsize$[-0.10,-0.05]$}} & {\boldmath $-1.526$ {\scriptsize$[-2.08,-0.35]$}} & $+0.000$ {\scriptsize$[+0.00,+0.00]$} \\
\addlinespace[2pt]
\multicolumn{4}{l}{(d) nested structure$+\rho$ $-$ policy (+calib.)}\\
\quad frozen ($n{=}105$) & {\boldmath $-0.093$ {\scriptsize$[-0.14,-0.04]$}} & $+0.004$ {\scriptsize$[-0.26,+0.16]$} & $-0.048$ {\scriptsize$[-0.14,+0.00]$} \\
\quad mixture ($n{=}523$) & {\boldmath $-0.330$ {\scriptsize$[-0.40,-0.28]$}} & {\boldmath $-2.355$ {\scriptsize$[-2.67,-0.86]$}} & $+0.000$ {\scriptsize$[+0.00,+0.00]$} \\
\midrule
\multicolumn{4}{l}{\emph{uniform jump}}\\
\midrule
\multicolumn{4}{l}{(a) nested structure-only $-$ policy}\\
\quad frozen ($n{=}105$) & $-0.003$ {\scriptsize$[-0.06,+0.06]$} & $+1.115$ {\scriptsize$[-0.21,+1.59]$} & $-0.047$ {\scriptsize$[-0.11,+0.00]$} \\
\quad churning ($n{=}419$) & $-0.009$ {\scriptsize$[-0.06,+0.07]$} & {\boldmath $+2.306$ {\scriptsize$[+0.17,+3.17]$}} & $-0.063$ {\scriptsize$[-0.15,+0.01]$} \\
\quad mixture ($n{=}524$) & {\boldmath $+0.104$ {\scriptsize$[+0.02,+0.20]$}} & {\boldmath $+4.166$ {\scriptsize$[+2.04,+5.46]$}} & {\boldmath $-0.222$ {\scriptsize$[-0.33,-0.12]$}} \\
\addlinespace[2pt]
\multicolumn{4}{l}{(b) $\rho$-only tree $-$ nested structure-only}\\
\quad frozen ($n{=}105$) & $-0.036$ {\scriptsize$[-0.09,+0.00]$} & $-1.306$ {\scriptsize$[-1.78,+0.00]$} & $+0.047$ {\scriptsize$[+0.00,+0.11]$} \\
\quad churning ($n{=}419$) & $-0.005$ {\scriptsize$[-0.08,+0.05]$} & $-1.976$ {\scriptsize$[-3.11,+0.00]$} & $-0.058$ {\scriptsize$[-0.15,+0.03]$} \\
\quad mixture ($n{=}524$) & {\boldmath $-0.123$ {\scriptsize$[-0.22,-0.04]$}} & {\boldmath $-3.836$ {\scriptsize$[-5.40,-1.55]$}} & $+0.125$ {\scriptsize$[-0.01,+0.26]$} \\
\addlinespace[2pt]
\multicolumn{4}{l}{(c) nested structure$+\rho$ $-$ policy}\\
\quad frozen ($n{=}105$) & $-0.003$ {\scriptsize$[-0.06,+0.06]$} & $+1.115$ {\scriptsize$[-0.21,+1.59]$} & $-0.047$ {\scriptsize$[-0.11,+0.00]$} \\
\quad churning ($n{=}419$) & {\boldmath $-0.060$ {\scriptsize$[-0.08,-0.04]$}} & $-0.168$ {\scriptsize$[-0.20,+0.02]$} & $+0.003$ {\scriptsize$[-0.06,+0.06]$} \\
\quad mixture ($n{=}524$) & {\boldmath $-0.044$ {\scriptsize$[-0.07,-0.02]$}} & $+1.338$ {\scriptsize$[-0.17,+2.04]$} & $-0.005$ {\scriptsize$[-0.07,+0.05]$} \\
\addlinespace[2pt]
\multicolumn{4}{l}{(d) nested structure$+\rho$ $-$ policy (+calib.)}\\
\quad frozen ($n{=}105$) & {\boldmath $-0.071$ {\scriptsize$[-0.12,-0.01]$}} & $+1.064$ {\scriptsize$[-0.27,+1.54]$} & $-0.047$ {\scriptsize$[-0.11,+0.00]$} \\
\quad churning ($n{=}419$) & {\boldmath $-0.304$ {\scriptsize$[-0.39,-0.23]$}} & {\boldmath $-0.981$ {\scriptsize$[-1.42,-0.46]$}} & $+0.003$ {\scriptsize$[-0.06,+0.06]$} \\
\quad mixture ($n{=}524$) & {\boldmath $-0.250$ {\scriptsize$[-0.31,-0.19]$}} & $+0.526$ {\scriptsize$[-1.42,+1.44]$} & $-0.005$ {\scriptsize$[-0.07,+0.05]$} \\
\bottomrule
\end{tabular}

\end{table}

\paragraph{Feature cost.} Of the six features, only $\mathrm{dens}$ is not free:
measured over 21 operators and five seeds it costs 7.5 column-exact applies,
because its implementation counts over every nonzero of $A$; with that count
hoisted out, the part that depends on the dirty set costs $0.16$ of an apply.
The probe's calibration excess is 115 applies under drift and 121 under jump, so
the probe costs $15$--$760\times$ more either way. Refitting the nested procedure
on strictly free metadata plus $\rho$ still beats the policy in all three scopes:
$-0.078$ $[-0.101, -0.054]$, $-0.060$ $[-0.095, -0.026]$, $-0.039$ $[-0.069,
-0.007]$. The win is not from $\rho$ alone --- $\rho$-only ties the constant arm
on uniform-jump churning cells, $1.0705$ against $1.0705$ --- and it is not from
$\mathrm{dens}$ either. It is from the free structural metadata and $\rho$
jointly.

\paragraph{Accuracy sits partly under the measurement floor.} $\rho{=}0$ is the
identity for both motion models, so the two $\rho{=}0$ blocks are the same
workload measured twice, and their disagreement is a measurement floor: median
$1.041$, p90 $1.117$, max $1.553$, with 18 of 105 over $1.10$. Under a flat
$10\%$ parity band, 53 of 419 uniform-jump churning cells are unresolvable, and
they carry $87\%$ of the $0.888$-versus-$0.833$ accuracy gap between feature sets
and $100\%$ of the offline-minus-policy accuracy gap. Regret is a different
matter and survives all three floor definitions: restricted to resolvable cells
the contrast keeps its sign and its interval. This is why accuracy is reported
in Table~\ref{tab:predictors} and is not used as an argument.

\begin{figure}[tbp]
  \centering
  \includegraphics[width=0.62\linewidth]{fig_regret_cdf}
  \caption{Where each strategy's cost lands, cell by cell, on the workload
  mixture. The geomean and the tail disagree about which method to prefer, which
  is why both are reported.}
  \label{fig:cdf}
\end{figure}

\paragraph{What the policy still has.} On the uniform-jump workload mixture the
best structure-only model's worst cell is $6.84\times$ against the policy's
$1.91\times$ in steady state, and the paired bootstrap on worst-case regret puts
the nested structure-only procedure $+4.17$ $[+2.04, +5.46]$ above the policy.
Selection here is variance control rather than mean improvement, and a benchmark
in this area should report the worst cell beside the geomean because the two
disagree about which method to prefer (Figure~\ref{fig:cdf}). The
structure-plus-$\rho$ model's worst cell is $2.03\times$, so most of that tail
protection is available offline as well. The policy also needs no training set
and no assumption that deployment operators resemble the operators a model was
fitted on; this sweep cannot price that advantage, because every model here is
scored leave-one-operator-out on the same 21 operators, all of them power flow,
circuit simulation or structural FEM.

\section{Amortization of the calibration cost}
\label{sec:amortize}

The probe is a one-time cost, so whether it amortizes is a question about
session length rather than about operator size, and the harness answers it
exactly. The reported per-step cost is the mean over all 400 steps and the
steady-state cost is the mean after the 32-step probe window, so the one-time
excess is exactly $400\,(c_{\text{full}} - c_{\text{steady}})$ per cell and an
$S$-step session costs $c_{\text{steady}} + \text{excess}/S$ per step. That is an
identity, not a fit.

\begin{figure}[tbp]
  \centering
  \includegraphics[width=0.62\linewidth]{fig_amortization}
  \caption{Fraction of the hindsight ceiling held, against session length, for
  the three scopes that carry a live decision: the local-drift workload mixture,
  uniform-jump churning cells, and the uniform-jump workload mixture. The
  sweep's own length is marked, and the three curves pass through
  $0.752/0.753/0.781$ there. Of the pair \S\ref{sec:select} quotes at 400 steps,
  $71.8\%$ and $75.3\%$, only the second is on this plot --- it is the middle
  curve, uniform-jump churning cells, at $0.753$. The $71.8\%$ is local-drift
  churning cells, which is not one of these three scopes and has no curve here.}
  \label{fig:amort}
\end{figure}

Held fraction of the ceiling is $1/\text{geomean(regret)}$, and over the three
scopes it rises from $0.27$--$0.32$ at $S{=}32$ through $0.75$--$0.78$ at
$S{=}400$ to $0.87$--$0.89$ at $S{=}1600$ and an asymptote of $0.92$--$0.93$
(Figure~\ref{fig:amort}). Contrasts (c) and (d) in
Table~\ref{tab:contrasts} are the same quantity at $S{=}\infty$ and $S{=}400$:
on the local-drift mixture the difference runs $-0.330$ at 400 steps, $-0.208$
at 800, $-0.145$ at 1600, $-0.096$ at 6400 and $-0.079$ in the limit. The
policy's long-session number is therefore $92$--$93\%$, against the
$0.75$--$0.78$ a 400-step harness reports.

At $92$--$93\%$ it is still behind. Against the nested procedure the policy never
overtakes, at any session length, under any of the three ways of charging for
$\mathrm{dens}$. Against the pre-registered logistic regression it overtakes at
$S \approx 1000$--$1300$ on the two mixtures, with the interval covering zero
from about $S{=}800$, and never on uniform-jump churning cells.

The curve extrapolates a single regime. Both motion models change regime at most
once in this benchmark, so nothing here prices a workload that alternates: each
real regime change costs the policy another probe and costs the offline model
only a feature recompute. Which of those is worse is untested and is the single
most likely thing to change the answer.

\section{A merged-but-patchable schedule}
\label{sec:patch}

Sections~\ref{sec:families} and~\ref{sec:merge} isolate two properties that
would be valuable together. A merged schedule touches $2.9\times$ fewer distinct
output rows on a frozen or locally drifting set; per-tile compiled segments are
what make re-preparation cheap when the set moves. We built and measured a third
implementation that has both: merged rows laid out with per-row slack so that a
single column-tile's contributions can be spliced in or out of a live schedule
without re-emitting it.

\paragraph{Layout.}
The merge survives being laid out that way exactly: the merged-and-patchable
schedule touches the identical number of distinct output rows as the flat merged
schedule on \textbf{1050 of 1050} cells of the 21-operator sweep, and on 165 of
165 of the five-operator timing files. That is an integer identity
rather than a timing, so the measurement floor cannot reach it. The expected
failure mode was that slack would dilute the schedule and
cost the kernel its locality. It does not: at occupancy ratios up to $5.95\times$
(drift) and $7.52\times$ (jump) the patchable kernel runs at $0.976\times$ the
packed one, a $2.4\%$ cost, geomean over 60 churning cells per model.

\paragraph{Incremental build.}
As a whole lane the merged-and-patchable strategy runs at $0.81\times$ the
per-tile-compiled one under local drift and $0.66\times$ under uniform jump; it
is ahead in 17 of 60 cells and 4 of 60. Splicing is the cheaper way to move a
merged schedule: one splice costs $1.11$--$105\times$ less than the from-scratch
merged rebuild it replaces, in all nine cells measured. It is not competitive
with never having merged (Figure~\ref{fig:splice}). Entering the patchable form
costs $2.7$--$3.7\times$ a flat merged build, so against a merged schedule that
rebuilds every step the crossover runs from $1.8$ to $27.3$ steps.

\begin{figure}[tbp]
  \centering
  \includegraphics[width=\linewidth]{fig_splice_events}
  \caption{Placement events and their price. The event-count gap reaches
  $50.5\times$ and the per-event ratio runs the other way, so neither number
  alone describes the strategy.}
  \label{fig:splice}
\end{figure}

\paragraph{Mechanism.}
A tile's contributions are context-free: they can be derived once and cached
against the tile, which is what the per-tile-compiled strategy does. What is not
context-free is \emph{placement} --- where a merged row's span sits, how much
slack it has, what its entry offsets are --- because all three depend on which
other tiles are currently present. The event counts follow. At the widest cell
the patchable strategy performs $55.6$ placement events per step where the
per-tile one compiles $1.1$, and because $|D|$ and the tile count are stationary
(at $194.1$), those events are about $27.8$ arrivals and $27.8$ departures.

The decomposition is over event counts and the conclusion a reader wants is
about cost, and we did not measure the two event classes separately, so what
follows is bounded rather than settled. Departures are irreducible for a merged
layout and free for a list of segments. Arrivals are not removable either, since
an arriving tile still has to be placed; what is cacheable is the
\emph{derivation} of an arriving tile's contributions, by exactly the mechanism
the per-tile strategy already has and which the implementation measured here
does not. A cache would therefore make arrivals cheaper without making them go
away, and the residual gap would still be $27.8$ placement events against $1.1$
compiles. Quoting the $50.5\times$ whole overstates the barrier; calling half of
it an unbuilt cache understates it. What this section settles is its layout
half, the row-count identity and the occupancy measurement, rather than its
mechanism half.

\paragraph{As a third arm.}
The same implementation was measured as a candidate third arm across all 21
operators and both motion models. It is at parity with the shipped
block-scheduled arm only on a frozen set ($1.12\times$ / $1.11\times$) and at the
lowest churn rate under local drift ($1.05\times$), and falls to $0.43\times$ and
$0.20\times$ by $\rho{=}0.25$. A three-arm hindsight oracle would beat the
two-arm one by at most $8.3\%$, at local drift $\rho{=}0.002$, and by $0.2\%$ or
less at every uniform-jump rate (Figure~\ref{fig:third}). Since the online
policy probes every arm, a third arm costs a third probe on every re-decision to
capture that, against a policy whose calibration cost is already what limits it
(Appendix~\ref{sec:amortize}).

\begin{figure}[tbp]
  \centering
  \includegraphics[width=0.86\linewidth]{fig_third_arm}
  \caption{The third implementation against the shipped block-scheduled arm, and
  what a three-arm hindsight oracle would have been worth over the two-arm one.}
  \label{fig:third}
\end{figure}

% otherwise this heading lands alone at the foot of the previous page
\clearpage
\section{Per-operator results}

% [!ht] rather than [tbp]: this is the only content in its appendix, and a
% top-of-page float put the table above its own section heading.
\begin{table}[!ht]
  \centering
  \caption{All 21 operators under both motion models. Values are the better
  arm's margin over the column-exact baseline, median across five dirty-set
  seeds. Bold marks cells where the block-scheduled arm is the better one.}
  \label{tab:operators}
  \scriptsize
  % eleven columns; the stock 6pt gutter put it 13.6pt past the text block
  \setlength{\tabcolsep}{4pt}
  \begin{tabular}{llrrrlccccc}
\toprule
operator & domain & $N$ & nnz & dens & motion & \multicolumn{5}{c}{best arm $\div$ column-exact, by $\rho$}\\
\cmidrule(lr){7-11}
 & & & & & & $0$ & $0.002$ & $0.01$ & $0.05$ & $0.25$ \\
\midrule
TSOPF\_RS\_b39\_c30 & power flow & 60,098 & 1,079,986 & 1.00 & local drift & \textbf{1.00} & \textbf{33.9} & \textbf{34.9} & \textbf{33.6} & \textbf{33.4} \\
 &  &  &  &  & uniform jump & \textbf{1.01} & \textbf{16.7} & \textbf{9.18} & \textbf{3.81} & \textbf{1.60} \\
TSOPF\_RS\_b162\_c3 & power flow & 15,374 & 610,299 & 1.00 & local drift & 1.00 & \textbf{31.9} & \textbf{30.9} & \textbf{28.2} & \textbf{29.0} \\
 &  &  &  &  & uniform jump & 1.00 & \textbf{14.3} & \textbf{7.87} & \textbf{4.18} & \textbf{2.57} \\
TSOPF\_RS\_b300\_c1 & power flow & 14,538 & 1,474,325 & 1.00 & local drift & 1.00 & \textbf{24.0} & \textbf{24.0} & \textbf{21.8} & \textbf{29.1} \\
 &  &  &  &  & uniform jump & 1.00 & \textbf{16.3} & \textbf{9.54} & \textbf{5.49} & \textbf{3.82} \\
nasa2910 & structural FEM & 2,910 & 174,296 & 0.84 & local drift & 1.00 & \textbf{14.2} & \textbf{12.2} & \textbf{8.57} & \textbf{8.70} \\
 &  &  &  &  & uniform jump & 1.00 & \textbf{5.37} & \textbf{4.29} & \textbf{3.85} & \textbf{3.71} \\
raefsky3 & fluid/struct. & 21,200 & 1,488,768 & 1.00 & local drift & 1.00 & \textbf{13.6} & \textbf{11.6} & \textbf{9.42} & \textbf{9.10} \\
 &  &  &  &  & uniform jump & 1.00 & \textbf{4.91} & \textbf{3.65} & \textbf{3.09} & \textbf{2.75} \\
case9 & power flow & 14,454 & 147,972 & 0.76 & local drift & 1.00 & \textbf{12.4} & \textbf{12.4} & \textbf{12.1} & \textbf{12.1} \\
 &  &  &  &  & uniform jump & 1.00 & \textbf{4.31} & \textbf{3.05} & \textbf{1.67} & \textbf{1.15} \\
msc10848 & structural FEM & 10,848 & 1,229,776 & 0.68 & local drift & 1.00 & \textbf{10.5} & \textbf{9.83} & \textbf{8.23} & \textbf{8.42} \\
 &  &  &  &  & uniform jump & 1.00 & \textbf{5.28} & \textbf{4.18} & \textbf{3.75} & \textbf{3.50} \\
ct20stif & structural FEM & 52,329 & 2,600,295 & 0.75 & local drift & 1.00 & \textbf{9.73} & \textbf{8.59} & \textbf{6.22} & \textbf{5.90} \\
 &  &  &  &  & uniform jump & 1.00 & \textbf{3.11} & \textbf{2.30} & \textbf{1.95} & \textbf{1.75} \\
s3rmt3m3 & shell FEM & 5,357 & 207,123 & 0.82 & local drift & 1.00 & \textbf{9.57} & \textbf{8.33} & \textbf{5.50} & \textbf{5.08} \\
 &  &  &  &  & uniform jump & 1.00 & \textbf{3.95} & \textbf{3.18} & \textbf{2.82} & \textbf{2.65} \\
memplus & circuit sim & 17,758 & 99,147 & 0.82 & local drift & 1.00 & \textbf{9.49} & \textbf{8.87} & \textbf{7.63} & \textbf{6.64} \\
 &  &  &  &  & uniform jump & 1.00 & \textbf{4.41} & \textbf{2.25} & \textbf{1.15} & 1.00 \\
olafu & structural FEM & 16,146 & 1,015,156 & 0.63 & local drift & 1.00 & \textbf{9.13} & \textbf{8.69} & \textbf{6.74} & \textbf{7.66} \\
 &  &  &  &  & uniform jump & 1.00 & \textbf{3.64} & \textbf{2.99} & \textbf{2.71} & \textbf{2.48} \\
bcsstk25 & structural FEM & 15,439 & 252,241 & 0.62 & local drift & 1.00 & \textbf{5.95} & \textbf{4.66} & \textbf{3.28} & \textbf{3.07} \\
 &  &  &  &  & uniform jump & 1.00 & \textbf{1.81} & \textbf{1.35} & \textbf{1.14} & \textbf{1.05} \\
circuit\_3 & circuit sim & 12,127 & 48,137 & 0.80 & local drift & 1.00 & \textbf{5.46} & \textbf{5.56} & \textbf{4.28} & \textbf{4.31} \\
 &  &  &  &  & uniform jump & 1.00 & \textbf{1.59} & 1.00 & 1.00 & 1.00 \\
bcsstk18 & structural FEM & 11,948 & 149,090 & 0.57 & local drift & 1.00 & \textbf{5.00} & \textbf{3.93} & \textbf{2.73} & \textbf{2.62} \\
 &  &  &  &  & uniform jump & 1.00 & \textbf{1.66} & \textbf{1.20} & \textbf{1.02} & 1.00 \\
add32 & circuit sim & 4,960 & 19,848 & 0.66 & local drift & 1.00 & \textbf{3.91} & \textbf{3.50} & \textbf{2.21} & \textbf{1.79} \\
 &  &  &  &  & uniform jump & 1.00 & \textbf{1.30} & \textbf{1.05} & 1.00 & 1.00 \\
epb2 & plate/beam & 25,228 & 175,027 & 0.69 & local drift & 1.00 & \textbf{3.59} & \textbf{2.92} & \textbf{2.03} & \textbf{1.98} \\
 &  &  &  &  & uniform jump & 1.00 & \textbf{1.21} & 1.00 & 1.00 & 1.00 \\
rajat03 & circuit sim & 7,602 & 32,653 & 0.45 & local drift & 1.00 & \textbf{2.88} & \textbf{2.16} & \textbf{1.51} & \textbf{1.28} \\
 &  &  &  &  & uniform jump & 1.00 & \textbf{1.10} & 1.00 & 1.00 & 1.00 \\
wang3 & device sim & 26,064 & 177,168 & 0.37 & local drift & 1.00 & \textbf{2.32} & \textbf{1.77} & \textbf{1.33} & \textbf{1.16} \\
 &  &  &  &  & uniform jump & 1.00 & \textbf{1.08} & 1.00 & 1.00 & 1.00 \\
scircuit & circuit sim & 170,998 & 958,936 & 0.48 & local drift & 1.00 & \textbf{2.24} & \textbf{1.84} & \textbf{1.57} & \textbf{1.43} \\
 &  &  &  &  & uniform jump & 1.00 & 1.00 & 1.00 & 1.00 & 1.00 \\
bcspwr09 & power network & 1,723 & 6,511 & 0.26 & local drift & 1.00 & \textbf{1.97} & \textbf{1.64} & \textbf{1.55} & \textbf{1.52} \\
 &  &  &  &  & uniform jump & 1.00 & \textbf{1.73} & \textbf{1.42} & \textbf{1.26} & \textbf{1.20} \\
bcspwr10 & power network & 5,300 & 21,842 & 0.18 & local drift & 1.00 & \textbf{1.32} & \textbf{1.26} & \textbf{1.22} & \textbf{1.15} \\
 &  &  &  &  & uniform jump & 1.00 & \textbf{1.20} & \textbf{1.06} & 1.00 & 1.00 \\
\bottomrule
\end{tabular}

\end{table}
\end{document}
